% !TEX root = iclr2027_conference.tex
% Fragment: % !TEX root = iclr2027_conference.tex
% Fragment: % !TEX root = iclr2027_conference.tex
% Fragment: % !TEX root = iclr2027_conference.tex
% Fragment: \input{pdcm_theory} from the main document.
% Requires in the preamble of the including file:
%   \usepackage{amsthm,mathtools}
%   \newtheorem{assumption}{Assumption}
%   \newtheorem{lemma}{Lemma}
%   \newtheorem{theorem}{Theorem}
%   \newtheorem{corollary}{Corollary}
%   \newtheorem{remark}{Remark}
%   \newcommand{\DeltaK}{\Delta_K}
%   \newcommand{\one}{\mathbf{1}}
%   \newcommand{\pos}[1]{\left[#1\right]_+}
%   \newcommand{\ip}[2]{\left\langle #1,#2\right\rangle}
%   \newcommand{\norm}[1]{\left\lVert #1\right\rVert}
% \R and \KL come from math_commands.tex (lines 480, 488).

\section{Theory}
\label{app:theory}

This appendix gives a self-contained analysis of the full-Jacobian version of PDCM used in Algorithm~\ref{alg:pdcm_jacobian}. The analysis isolates the online allocation mechanism. At each round it treats the allocation-response Jacobian as available exactly.

\subsection{Problem formulation and idealized full-Jacobian PDCM}
\label{app:theory_setup}

We consider $K\ge 2$ tasks and the probability simplex
\begin{equation}
    \Delta_K
    \coloneqq
    \left\{p\in\mathbb{R}_+^K:\mathbf{1}^\top p=1\right\},
    \qquad
    u\coloneqq \frac{1}{K}\mathbf{1}.
\end{equation}
At round $t$, let
\begin{equation}
    U_t(p)
    \coloneqq
    \begin{bmatrix}
        U_{t,1}(p) & \cdots & U_{t,K}(p)
    \end{bmatrix}^{\!\top}\in\mathbb{R}^K
\end{equation}
denote the vector of expected one-step policy improvements that would result from applying the round-$t$ update under allocation $p\in\Delta_K$. The aggregate improvement is
\begin{equation}
    F_t(p)\coloneqq \mathbf{1}^\top U_t(p)=\sum_{k=1}^K U_{t,k}(p).
    \label{eq:app_aggregate_utility}
\end{equation}
Because the policy state changes with training, the functions $U_t$ may be nonstationary in $t$.

We use the following structural condition. It is stronger than the reference-direction inequality needed only for the safety certificate, but it is the condition required to compare PDCM with an arbitrary fixed allocation in the regret statement.

\begin{assumption}[Concave allocation response]
\label{ass:app_concavity}
For every round $t$ and task $k$, $U_{t,k}$ is differentiable and concave on $\Delta_K$. Equivalently, for every $p,q\in\Delta_K$ and $\alpha\in[0,1]$,
\begin{equation}
    U_{t,k}((1-\alpha)p+\alpha q)
    \ge
    (1-\alpha)U_{t,k}(p)+\alpha U_{t,k}(q).
    \label{eq:app_line_concavity}
\end{equation}
\end{assumption}

Assumption~\ref{ass:app_concavity} implies the usual first-order upper bound for a concave function.

\begin{lemma}[Tangent upper bound]
\label{lem:app_tangent}
Under Assumption~\ref{ass:app_concavity}, for every $t,k$ and $p,q\in\Delta_K$,
\begin{equation}
    U_{t,k}(q)-U_{t,k}(p)
    \le
    \left\langle \nabla U_{t,k}(p),q-p\right\rangle.
    \label{eq:app_tangent_bound}
\end{equation}
\end{lemma}

\begin{proof}
For $\alpha\in(0,1]$, concavity gives
\begin{equation}
    U_{t,k}\!\left(p+\alpha(q-p)\right)-U_{t,k}(p)
    \ge
    \alpha\bigl(U_{t,k}(q)-U_{t,k}(p)\bigr).
\end{equation}
Divide by $\alpha$ and let $\alpha\downarrow 0$. Differentiability yields Equation~\eqref{eq:app_tangent_bound}.
\end{proof}

At the played allocation $p_t$, define the allocation-response Jacobian
\begin{equation}
    G_t
    \coloneqq
    \nabla_p U_t(p_t)
    =
    \begin{bmatrix}
        g_{t,1}^{\top}\\
        \vdots\\
        g_{t,K}^{\top}
    \end{bmatrix}
    \in\mathbb{R}^{K\times K},
    \qquad
    g_{t,k}\coloneqq \nabla_p U_{t,k}(p_t).
    \label{eq:app_jacobian}
\end{equation}
Thus $[G_t]_{kj}=\partial U_{t,k}(p_t)/\partial p_j$ measures the local effect of allocating additional mass to training task $j$ on the improvement of evaluation task $k$. In particular,
\begin{equation}
    \nabla F_t(p_t)=G_t^\top\mathbf{1}.
    \label{eq:app_aggregate_gradient}
\end{equation}

For task $k$, let $\epsilon_{t,k}\ge 0$ denote the lag budget assigned to round $t$. For a total cumulative budget $\epsilon_k$, one may choose any schedule satisfying $\sum_{t=1}^T\epsilon_{t,k}=\epsilon_k$, for example $\epsilon_{t,k}=\epsilon_k/T$. Define the true one-step counterfactual lag residual
\begin{equation}
    c_{t,k}(p_t)
    \coloneqq
    U_{t,k}(u)-U_{t,k}(p_t)-\epsilon_{t,k}.
    \label{eq:app_true_residual}
\end{equation}
Applying Lemma~\ref{lem:app_tangent} with $p=p_t$ and $q=u$ gives
\begin{equation}
    c_{t,k}(p_t)
    \le
    \left\langle g_{t,k},u-p_t\right\rangle-\epsilon_{t,k}.
    \label{eq:app_safety_certificate}
\end{equation}
Therefore the first-order quantity on the right is a conservative certificate of the unobserved uniform-allocation lag.

For the online analysis, extend this certificate affinely to an arbitrary comparator $p\in\Delta_K$:
\begin{equation}
    h_t(p)
    \coloneqq
    G_t(u-p)-\epsilon_t,
    \qquad
    \epsilon_t
    \coloneqq
    (\epsilon_{t,1},\ldots,\epsilon_{t,K})^\top.
    \label{eq:app_affine_residual}
\end{equation}
Equation~\eqref{eq:app_safety_certificate} is equivalently
\begin{equation}
    c_t(p_t)\le h_t(p_t)
    \qquad\text{coordinatewise}.
    \label{eq:app_true_vs_linearized}
\end{equation}
For horizon $T$, define the fixed first-order-safe comparator set
\begin{equation}
    \mathcal{P}_T
    \coloneqq
    \left\{
        p\in\Delta_K:
        h_t(p)\le 0\ \text{coordinatewise for every }t\le T
    \right\}.
    \label{eq:app_comparator_set}
\end{equation}
This set is always nonempty because $h_t(u)=-\epsilon_t\le 0$.

Let $\lambda_t\in\mathbb{R}_+^K$ be the task-wise dual variables. Dropping terms independent of $p$ from the linearized round-$t$ Lagrangian gives the primal score
\begin{equation}
    s_t
    \coloneqq
    G_t^\top(\mathbf{1}+\lambda_t).
    \label{eq:app_score}
\end{equation}
The idealized full-Jacobian PDCM recursion is
\begin{align}
    p_{t+1}
    &\coloneqq
    \arg\max_{p\in\Delta_K}
    \left\{
        \eta_p\langle s_t,p\rangle
        -D_{\mathrm{KL}}(p\|p_t)
    \right\},
    \label{eq:app_primal_update}\\
    \lambda_{t+1}
    &\coloneqq
    \left[
        \lambda_t+\eta_\lambda h_t(p_t)
    \right]_+,
    \label{eq:app_dual_update}
\end{align}
with $p_1=u$ and $\lambda_1=0$. Here $[\cdot]_+$ denotes coordinatewise projection onto $\mathbb{R}_+^K$. The primal step has the closed form
\begin{equation}
    p_{t+1,j}
    =
    \frac{
        p_{t,j}\exp(\eta_p s_{t,j})
    }{
        \sum_{\ell=1}^K p_{t,\ell}\exp(\eta_p s_{t,\ell})
    }.
    \label{eq:app_exp_update}
\end{equation}

We impose only boundedness conditions needed by the finite-horizon analysis.

\begin{assumption}[Bounded Jacobian and lag budget]
\label{ass:app_bounded}
There exist finite constants $L,E$ such that for every round $t$,
\begin{equation}
    \|G_t\|_1
    \coloneqq
    \max_{j\in[K]}\sum_{k=1}^K |[G_t]_{kj}|
    \le L,
    \qquad
    0\le \epsilon_{t,k}\le E.
    \label{eq:app_boundedness}
\end{equation}
\end{assumption}

For a horizon $T$, define the realized dual magnitude
\begin{equation}
    \Lambda_T
    \coloneqq
    \max_{1\le t\le T+1}\|\lambda_t\|_\infty,
    \qquad
    H\coloneqq 2L+E.
    \label{eq:app_constants}
\end{equation}
Then Assumption~\ref{ass:app_bounded} implies
\begin{equation}
    \|s_t\|_\infty\le L(1+\Lambda_T),
    \qquad
    \|h_t(p_t)\|_\infty\le H.
    \label{eq:app_basic_bounds}
\end{equation}
The second inequality follows from $\|u-p_t\|_1\le 2$ and $|[G_t]_{kj}|\le L$.

\subsection{Constrained no-regret guarantee}
\label{app:theory_regret}

For a fixed comparator $p\in\Delta_K$, define cumulative utility regret
\begin{equation}
    R_T(p)
    \coloneqq
    \sum_{t=1}^T\bigl(F_t(p)-F_t(p_t)\bigr).
    \label{eq:app_regret}
\end{equation}
The next theorem gives the finite-horizon bound from which the $O(\sqrt{T})$ statement in the main text follows.

\begin{theorem}[Constrained no-regret of full-Jacobian PDCM]
\label{thm:app_no_regret}
Suppose Assumptions~\ref{ass:app_concavity}--\ref{ass:app_bounded} hold and $(p_t,\lambda_t)$ is generated by Equations~\eqref{eq:app_primal_update}--\eqref{eq:app_dual_update}. Then for every horizon $T\ge 1$, every $p\in\mathcal{P}_T$, and every $\eta_p,\eta_\lambda>0$,
\begin{equation}
    R_T(p)
    \le
    \frac{\log K}{\eta_p}
    +
    \frac{\eta_p T}{2}L^2(1+\Lambda_T)^2
    +
    \frac{\eta_\lambda T}{2}KH^2.
    \label{eq:app_finite_regret_bound}
\end{equation}
Moreover, for every task $k$,
\begin{equation}
    \left[
        \sum_{t=1}^T c_{t,k}(p_t)
    \right]_+
    \le
    \left[
        \sum_{t=1}^T h_{t,k}(p_t)
    \right]_+
    \le
    \frac{\Lambda_T}{\eta_\lambda}.
    \label{eq:app_constraint_bound}
\end{equation}
If the dual sequence is stable, i.e., $\Lambda_T\le\Lambda<\infty$ for all $T$, and $\eta_p,\eta_\lambda=\Theta(T^{-1/2})$, then
\begin{equation}
    R_T(p)=O(\sqrt{T}),
    \qquad
    \left[
        \sum_{t=1}^T c_{t,k}(p_t)
    \right]_+
    =O(\sqrt{T}),
    \label{eq:app_sqrt_rates}
\end{equation}
so both average regret and average excess counterfactual lag vanish as $O(T^{-1/2})$.
\end{theorem}

\begin{proof}
We first prove a one-step entropic mirror-ascent inequality. By optimality of Equation~\eqref{eq:app_primal_update} and the three-point identity for the negative-entropy Bregman divergence, for every $p\in\Delta_K$,
\begin{equation}
    \eta_p\langle s_t,p-p_{t+1}\rangle
    \le
    D_{\mathrm{KL}}(p\|p_t)
    -D_{\mathrm{KL}}(p\|p_{t+1})
    -D_{\mathrm{KL}}(p_{t+1}\|p_t).
    \label{eq:app_three_point}
\end{equation}
Adding $\eta_p\langle s_t,p_{t+1}-p_t\rangle$ to both sides and using H\"older's inequality together with Pinsker's inequality,
\begin{equation}
    \langle s_t,p_{t+1}-p_t\rangle
    \le
    \|s_t\|_\infty\|p_{t+1}-p_t\|_1,
    \qquad
    D_{\mathrm{KL}}(p_{t+1}\|p_t)
    \ge
    \frac{1}{2}\|p_{t+1}-p_t\|_1^2,
\end{equation}
we obtain
\begin{equation}
    \langle s_t,p-p_t\rangle
    \le
    \frac{
        D_{\mathrm{KL}}(p\|p_t)
        -D_{\mathrm{KL}}(p\|p_{t+1})
    }{\eta_p}
    +
    \frac{\eta_p}{2}\|s_t\|_\infty^2.
    \label{eq:app_mirror_one_step}
\end{equation}
Summing over $t$, using $p_1=u$, $D_{\mathrm{KL}}(p\|u)\le\log K$, and Equation~\eqref{eq:app_basic_bounds}, gives
\begin{equation}
    \sum_{t=1}^T\langle s_t,p-p_t\rangle
    \le
    \frac{\log K}{\eta_p}
    +
    \frac{\eta_p T}{2}L^2(1+\Lambda_T)^2.
    \label{eq:app_mirror_sum}
\end{equation}

We now relate the score to utility regret and the safety residuals. From Equation~\eqref{eq:app_affine_residual},
\begin{equation}
    h_t(p_t)-h_t(p)=G_t(p-p_t).
    \label{eq:app_h_difference}
\end{equation}
Therefore, using Equation~\eqref{eq:app_score},
\begin{align}
    \langle s_t,p-p_t\rangle
    &=
    \left\langle G_t^\top\mathbf{1},p-p_t\right\rangle
    +
    \left\langle \lambda_t,G_t(p-p_t)\right\rangle\\
    &=
    \left\langle G_t^\top\mathbf{1},p-p_t\right\rangle
    +
    \left\langle \lambda_t,h_t(p_t)-h_t(p)\right\rangle.
    \label{eq:app_score_decomposition}
\end{align}
For $p\in\mathcal{P}_T$, we have $h_t(p)\le0$ coordinatewise and $\lambda_t\ge0$, so
\begin{equation}
    \left\langle G_t^\top\mathbf{1},p-p_t\right\rangle
    +
    \langle\lambda_t,h_t(p_t)\rangle
    \le
    \langle s_t,p-p_t\rangle.
    \label{eq:app_safe_drop}
\end{equation}
Since each $U_{t,k}$ is concave, $F_t$ is concave, and Lemma~\ref{lem:app_tangent} implies
\begin{equation}
    F_t(p)-F_t(p_t)
    \le
    \langle\nabla F_t(p_t),p-p_t\rangle
    =
    \left\langle G_t^\top\mathbf{1},p-p_t\right\rangle.
    \label{eq:app_F_concavity}
\end{equation}
Combining Equations~\eqref{eq:app_mirror_sum}, \eqref{eq:app_safe_drop}, and \eqref{eq:app_F_concavity} yields
\begin{equation}
    R_T(p)
    +
    \sum_{t=1}^T\langle\lambda_t,h_t(p_t)\rangle
    \le
    \frac{\log K}{\eta_p}
    +
    \frac{\eta_p T}{2}L^2(1+\Lambda_T)^2.
    \label{eq:app_predual}
\end{equation}

It remains to control the dual term. Projection onto the nonnegative orthant is nonexpansive and fixes the origin, hence Equation~\eqref{eq:app_dual_update} gives
\begin{align}
    \|\lambda_{t+1}\|_2^2
    &\le
    \|\lambda_t+\eta_\lambda h_t(p_t)\|_2^2\\
    &=
    \|\lambda_t\|_2^2
    +2\eta_\lambda\langle\lambda_t,h_t(p_t)\rangle
    +\eta_\lambda^2\|h_t(p_t)\|_2^2.
    \label{eq:app_dual_square}
\end{align}
Rearranging and summing from $t=1$ to $T$, using $\lambda_1=0$ and $\|h_t(p_t)\|_2^2\le KH^2$, gives
\begin{equation}
    -\sum_{t=1}^T\langle\lambda_t,h_t(p_t)\rangle
    \le
    \frac{\eta_\lambda T}{2}KH^2.
    \label{eq:app_dual_inner_bound}
\end{equation}
Substituting Equation~\eqref{eq:app_dual_inner_bound} into Equation~\eqref{eq:app_predual} proves Equation~\eqref{eq:app_finite_regret_bound}.

For the constraint bound, the coordinatewise projection satisfies
\begin{equation}
    \lambda_{t+1,k}
    =
    \max\{0,\lambda_{t,k}+\eta_\lambda h_{t,k}(p_t)\}
    \ge
    \lambda_{t,k}+\eta_\lambda h_{t,k}(p_t).
\end{equation}
Telescoping and using $\lambda_{1,k}=0$ gives
\begin{equation}
    \eta_\lambda\sum_{t=1}^T h_{t,k}(p_t)
    \le
    \lambda_{T+1,k}
    \le
    \Lambda_T.
\end{equation}
Taking positive parts proves the second inequality in Equation~\eqref{eq:app_constraint_bound}. The first follows from Equation~\eqref{eq:app_true_vs_linearized} after summing over $t$. Finally, if $\Lambda_T\le\Lambda$ and both step sizes are of order $T^{-1/2}$, every term in Equations~\eqref{eq:app_finite_regret_bound} and \eqref{eq:app_constraint_bound} is $O(\sqrt{T})$.
\end{proof}

The preceding theorem immediately yields the comparator used in the main text.

\begin{corollary}[Best fixed feasible allocation]
\label{cor:app_best_fixed}
Let
\begin{equation}
    p_T^\star
    \in
    \arg\max_{p\in\mathcal{P}_T}
    \sum_{t=1}^T F_t(p).
    \label{eq:app_best_fixed}
\end{equation}
Under the dual-stability condition $\Lambda_T\le\Lambda$ and step sizes $\eta_p,\eta_\lambda=\Theta(T^{-1/2})$,
\begin{equation}
    \sum_{t=1}^T
    \left[F_t(p_T^\star)-F_t(p_t)\right]
    =O(\sqrt{T}),
    \label{eq:app_best_fixed_regret}
\end{equation}
and for every task $k$,
\begin{equation}
    \left[
        \sum_{t=1}^T
        \left(
            U_{t,k}(u)-U_{t,k}(p_t)-\epsilon_{t,k}
        \right)
    \right]_+
    =O(\sqrt{T}).
    \label{eq:app_best_fixed_constraint}
\end{equation}
\end{corollary}

\begin{corollary}[Uniform-mixture safeguard]
\label{cor:app_uniform}
Because $u\in\mathcal{P}_T$, Theorem~\ref{thm:app_no_regret} applies in particular to the uniform mixture. Under dual stability and $\eta_p,\eta_\lambda=\Theta(T^{-1/2})$,
\begin{equation}
    \sum_{t=1}^T F_t(u)
    -
    \sum_{t=1}^T F_t(p_t)
    =O(\sqrt{T}).
    \label{eq:app_uniform_regret}
\end{equation}
Hence the average aggregate-utility gap between PDCM and the static uniform mixture vanishes as $O(T^{-1/2})$. In addition,
\begin{equation}
    \frac{1}{T}\sum_{t=1}^T
    \left(U_{t,k}(u)-U_{t,k}(p_t)\right)
    \le
    \frac{1}{T}\sum_{t=1}^T\epsilon_{t,k}
    +O(T^{-1/2})
    \label{eq:app_uniform_lag}
\end{equation}
for every task $k$.
\end{corollary}

\begin{remark}[Dual stability]
\label{rem:app_dual_stability}
The finite-horizon bounds in Equations~\eqref{eq:app_finite_regret_bound}--\eqref{eq:app_constraint_bound} hold for the realized quantity $\Lambda_T$ and do not assume that the dual sequence is bounded in advance. The $O(\sqrt{T})$ rates require $\Lambda_T=O(1)$. This is a stability condition on the coupled allocation and training dynamics; it may be checked empirically, imposed as an assumption, or enforced by an upper projection on the dual variables when such a projection is inactive relative to the desired optimum.
\end{remark}

\subsection{Dual-weighted marginal-utility matching}
\label{app:theory_matching}

The no-regret result describes finite-horizon behavior. A complementary stationary characterization explains the allocation condition targeted by the primal-dual coupling. Let $U_k:\Delta_K\to\mathbb{R}$ be differentiable and concave, define
\begin{equation}
    U(p)
    \coloneqq
    \begin{bmatrix}
        U_1(p)&\cdots&U_K(p)
    \end{bmatrix}^{\!\top},
    \qquad
    G(p)\coloneqq\nabla_p U(p),
    \qquad
    F(p)\coloneqq\mathbf{1}^\top U(p),
\end{equation}
and consider the static constrained allocation problem
\begin{equation}
\begin{aligned}
    \max_{p\in\Delta_K}\quad
    &F(p)\\
    \text{subject to}\quad
    &U_k(u)-U_k(p)-\epsilon_k\le0,
    \qquad k\in[K].
\end{aligned}
\label{eq:app_stationary_problem}
\end{equation}
Since $F$ is concave and each constraint function $U_k(u)-U_k(p)-\epsilon_k$ is convex, Equation~\eqref{eq:app_stationary_problem} is a convex optimization problem after negating the objective.

\begin{theorem}[Dual-weighted marginal-utility matching]
\label{thm:app_matching}
Assume Equation~\eqref{eq:app_stationary_problem} satisfies Slater's condition, and let $(p^\star,\lambda^\star)$ be a primal-dual optimum. Define
\begin{equation}
    s^\star
    \coloneqq
    G(p^\star)^\top(\mathbf{1}+\lambda^\star).
    \label{eq:app_stationary_score}
\end{equation}
Then there exists $\mu^\star\in\mathbb{R}$ such that
\begin{equation}
    s_j^\star=\mu^\star
    \quad\text{for every }j\text{ with }p_j^\star>0,
    \qquad
    s_j^\star\le\mu^\star
    \quad\text{for every }j\text{ with }p_j^\star=0.
    \label{eq:app_matching_condition}
\end{equation}
The multipliers additionally satisfy
\begin{equation}
    \lambda_k^\star\ge0,
    \qquad
    U_k(u)-U_k(p^\star)-\epsilon_k\le0,
    \qquad
    \lambda_k^\star\bigl(U_k(u)-U_k(p^\star)-\epsilon_k\bigr)=0.
    \label{eq:app_complementarity}
\end{equation}
If $p_j^\star>0$ for every $j$, then $p^\star$ is a fixed point of the exponentiated-gradient primal map with $\lambda^\star$ if and only if $s^\star=\mu^\star\mathbf{1}$.
\end{theorem}

\begin{proof}
Rewrite Equation~\eqref{eq:app_stationary_problem} as the convex minimization problem
\begin{equation}
\begin{aligned}
    \min_{p\in\Delta_K}\quad
    &-F(p)\\
    \text{subject to}\quad
    &c_k(p)\coloneqq U_k(u)-U_k(p)-\epsilon_k\le0,
    \qquad k\in[K].
\end{aligned}
\end{equation}
Under Slater's condition, the Karush--Kuhn--Tucker conditions are necessary and sufficient. Introduce multipliers $\lambda_k\ge0$ for $c_k(p)\le0$, $\mu\in\mathbb{R}$ for $\mathbf{1}^\top p=1$, and $\alpha_j\ge0$ for $-p_j\le0$. The KKT Lagrangian is
\begin{equation}
    \mathcal{L}_{\mathrm{KKT}}(p,\lambda,\mu,\alpha)
    =
    -F(p)
    +\sum_{k=1}^K\lambda_k\bigl(U_k(u)-U_k(p)-\epsilon_k\bigr)
    +\mu(\mathbf{1}^\top p-1)
    -\alpha^\top p.
\end{equation}
Differentiating with respect to $p$ gives the stationarity condition
\begin{equation}
    -G(p^\star)^\top(\mathbf{1}+\lambda^\star)
    +\mu^\star\mathbf{1}
    -\alpha^\star
    =0.
    \label{eq:app_KKT_stationarity}
\end{equation}
Hence $s_j^\star=\mu^\star-\alpha_j^\star$. If $p_j^\star>0$, complementary slackness for $-p_j\le0$ gives $\alpha_j^\star=0$, so $s_j^\star=\mu^\star$. If $p_j^\star=0$, then $\alpha_j^\star\ge0$, so $s_j^\star\le\mu^\star$. The remaining KKT conditions give Equation~\eqref{eq:app_complementarity}.

For the fixed-point statement, suppose first that $p_j^\star>0$ for every $j$ and $s^\star=\mu^\star\mathbf{1}$. Then the exponentiated-gradient map gives
\begin{equation}
    \frac{
        p_j^\star\exp(\eta_p s_j^\star)
    }{
        \sum_{\ell=1}^K p_\ell^\star\exp(\eta_p s_\ell^\star)
    }
    =
    \frac{
        p_j^\star\exp(\eta_p\mu^\star)
    }{
        \exp(\eta_p\mu^\star)\sum_{\ell=1}^K p_\ell^\star
    }
    =p_j^\star.
\end{equation}
Conversely, if the map fixes an interior $p^\star$, division by $p_j^\star>0$ shows that $\exp(\eta_p s_j^\star)$ equals the same normalization constant for every $j$. Taking logarithms gives $s_j^\star=s_\ell^\star$ for all $j,\ell$, which is exactly $s^\star=\mu^\star\mathbf{1}$ for some scalar $\mu^\star$.
\end{proof}

Theorem~\ref{thm:app_matching} exposes the role of the full Jacobian. Expanding the stationary score coordinatewise,
\begin{equation}
    s_j^\star
    =
    \sum_{k=1}^K
    (1+\lambda_k^\star)
    \frac{\partial U_k(p^\star)}{\partial p_j}.
    \label{eq:app_cross_task_matching}
\end{equation}
Thus the dual variable of task $k$ amplifies the entire $k$th row of the allocation-response Jacobian, not only its diagonal entry. A training coordinate $j$ is valuable when increasing its allocation improves any evaluation task with high marginal utility, and its score receives additional weight when those tasks are under active lag pressure. At an interior stationary optimum, these dual-weighted cross-task marginal effects are equalized across all training coordinates.
 from the main document.
% Requires in the preamble of the including file:
%   \usepackage{amsthm,mathtools}
%   \newtheorem{assumption}{Assumption}
%   \newtheorem{lemma}{Lemma}
%   \newtheorem{theorem}{Theorem}
%   \newtheorem{corollary}{Corollary}
%   \newtheorem{remark}{Remark}
%   \newcommand{\DeltaK}{\Delta_K}
%   \newcommand{\one}{\mathbf{1}}
%   \newcommand{\pos}[1]{\left[#1\right]_+}
%   \newcommand{\ip}[2]{\left\langle #1,#2\right\rangle}
%   \newcommand{\norm}[1]{\left\lVert #1\right\rVert}
% \R and \KL come from math_commands.tex (lines 480, 488).

\section{Theory}
\label{app:theory}

This appendix gives a self-contained analysis of the full-Jacobian version of PDCM used in Algorithm~\ref{alg:pdcm_jacobian}. The analysis isolates the online allocation mechanism. At each round it treats the allocation-response Jacobian as available exactly.

\subsection{Problem formulation and idealized full-Jacobian PDCM}
\label{app:theory_setup}

We consider $K\ge 2$ tasks and the probability simplex
\begin{equation}
    \Delta_K
    \coloneqq
    \left\{p\in\mathbb{R}_+^K:\mathbf{1}^\top p=1\right\},
    \qquad
    u\coloneqq \frac{1}{K}\mathbf{1}.
\end{equation}
At round $t$, let
\begin{equation}
    U_t(p)
    \coloneqq
    \begin{bmatrix}
        U_{t,1}(p) & \cdots & U_{t,K}(p)
    \end{bmatrix}^{\!\top}\in\mathbb{R}^K
\end{equation}
denote the vector of expected one-step policy improvements that would result from applying the round-$t$ update under allocation $p\in\Delta_K$. The aggregate improvement is
\begin{equation}
    F_t(p)\coloneqq \mathbf{1}^\top U_t(p)=\sum_{k=1}^K U_{t,k}(p).
    \label{eq:app_aggregate_utility}
\end{equation}
Because the policy state changes with training, the functions $U_t$ may be nonstationary in $t$.

We use the following structural condition. It is stronger than the reference-direction inequality needed only for the safety certificate, but it is the condition required to compare PDCM with an arbitrary fixed allocation in the regret statement.

\begin{assumption}[Concave allocation response]
\label{ass:app_concavity}
For every round $t$ and task $k$, $U_{t,k}$ is differentiable and concave on $\Delta_K$. Equivalently, for every $p,q\in\Delta_K$ and $\alpha\in[0,1]$,
\begin{equation}
    U_{t,k}((1-\alpha)p+\alpha q)
    \ge
    (1-\alpha)U_{t,k}(p)+\alpha U_{t,k}(q).
    \label{eq:app_line_concavity}
\end{equation}
\end{assumption}

Assumption~\ref{ass:app_concavity} implies the usual first-order upper bound for a concave function.

\begin{lemma}[Tangent upper bound]
\label{lem:app_tangent}
Under Assumption~\ref{ass:app_concavity}, for every $t,k$ and $p,q\in\Delta_K$,
\begin{equation}
    U_{t,k}(q)-U_{t,k}(p)
    \le
    \left\langle \nabla U_{t,k}(p),q-p\right\rangle.
    \label{eq:app_tangent_bound}
\end{equation}
\end{lemma}

\begin{proof}
For $\alpha\in(0,1]$, concavity gives
\begin{equation}
    U_{t,k}\!\left(p+\alpha(q-p)\right)-U_{t,k}(p)
    \ge
    \alpha\bigl(U_{t,k}(q)-U_{t,k}(p)\bigr).
\end{equation}
Divide by $\alpha$ and let $\alpha\downarrow 0$. Differentiability yields Equation~\eqref{eq:app_tangent_bound}.
\end{proof}

At the played allocation $p_t$, define the allocation-response Jacobian
\begin{equation}
    G_t
    \coloneqq
    \nabla_p U_t(p_t)
    =
    \begin{bmatrix}
        g_{t,1}^{\top}\\
        \vdots\\
        g_{t,K}^{\top}
    \end{bmatrix}
    \in\mathbb{R}^{K\times K},
    \qquad
    g_{t,k}\coloneqq \nabla_p U_{t,k}(p_t).
    \label{eq:app_jacobian}
\end{equation}
Thus $[G_t]_{kj}=\partial U_{t,k}(p_t)/\partial p_j$ measures the local effect of allocating additional mass to training task $j$ on the improvement of evaluation task $k$. In particular,
\begin{equation}
    \nabla F_t(p_t)=G_t^\top\mathbf{1}.
    \label{eq:app_aggregate_gradient}
\end{equation}

For task $k$, let $\epsilon_{t,k}\ge 0$ denote the lag budget assigned to round $t$. For a total cumulative budget $\epsilon_k$, one may choose any schedule satisfying $\sum_{t=1}^T\epsilon_{t,k}=\epsilon_k$, for example $\epsilon_{t,k}=\epsilon_k/T$. Define the true one-step counterfactual lag residual
\begin{equation}
    c_{t,k}(p_t)
    \coloneqq
    U_{t,k}(u)-U_{t,k}(p_t)-\epsilon_{t,k}.
    \label{eq:app_true_residual}
\end{equation}
Applying Lemma~\ref{lem:app_tangent} with $p=p_t$ and $q=u$ gives
\begin{equation}
    c_{t,k}(p_t)
    \le
    \left\langle g_{t,k},u-p_t\right\rangle-\epsilon_{t,k}.
    \label{eq:app_safety_certificate}
\end{equation}
Therefore the first-order quantity on the right is a conservative certificate of the unobserved uniform-allocation lag.

For the online analysis, extend this certificate affinely to an arbitrary comparator $p\in\Delta_K$:
\begin{equation}
    h_t(p)
    \coloneqq
    G_t(u-p)-\epsilon_t,
    \qquad
    \epsilon_t
    \coloneqq
    (\epsilon_{t,1},\ldots,\epsilon_{t,K})^\top.
    \label{eq:app_affine_residual}
\end{equation}
Equation~\eqref{eq:app_safety_certificate} is equivalently
\begin{equation}
    c_t(p_t)\le h_t(p_t)
    \qquad\text{coordinatewise}.
    \label{eq:app_true_vs_linearized}
\end{equation}
For horizon $T$, define the fixed first-order-safe comparator set
\begin{equation}
    \mathcal{P}_T
    \coloneqq
    \left\{
        p\in\Delta_K:
        h_t(p)\le 0\ \text{coordinatewise for every }t\le T
    \right\}.
    \label{eq:app_comparator_set}
\end{equation}
This set is always nonempty because $h_t(u)=-\epsilon_t\le 0$.

Let $\lambda_t\in\mathbb{R}_+^K$ be the task-wise dual variables. Dropping terms independent of $p$ from the linearized round-$t$ Lagrangian gives the primal score
\begin{equation}
    s_t
    \coloneqq
    G_t^\top(\mathbf{1}+\lambda_t).
    \label{eq:app_score}
\end{equation}
The idealized full-Jacobian PDCM recursion is
\begin{align}
    p_{t+1}
    &\coloneqq
    \arg\max_{p\in\Delta_K}
    \left\{
        \eta_p\langle s_t,p\rangle
        -D_{\mathrm{KL}}(p\|p_t)
    \right\},
    \label{eq:app_primal_update}\\
    \lambda_{t+1}
    &\coloneqq
    \left[
        \lambda_t+\eta_\lambda h_t(p_t)
    \right]_+,
    \label{eq:app_dual_update}
\end{align}
with $p_1=u$ and $\lambda_1=0$. Here $[\cdot]_+$ denotes coordinatewise projection onto $\mathbb{R}_+^K$. The primal step has the closed form
\begin{equation}
    p_{t+1,j}
    =
    \frac{
        p_{t,j}\exp(\eta_p s_{t,j})
    }{
        \sum_{\ell=1}^K p_{t,\ell}\exp(\eta_p s_{t,\ell})
    }.
    \label{eq:app_exp_update}
\end{equation}

We impose only boundedness conditions needed by the finite-horizon analysis.

\begin{assumption}[Bounded Jacobian and lag budget]
\label{ass:app_bounded}
There exist finite constants $L,E$ such that for every round $t$,
\begin{equation}
    \|G_t\|_1
    \coloneqq
    \max_{j\in[K]}\sum_{k=1}^K |[G_t]_{kj}|
    \le L,
    \qquad
    0\le \epsilon_{t,k}\le E.
    \label{eq:app_boundedness}
\end{equation}
\end{assumption}

For a horizon $T$, define the realized dual magnitude
\begin{equation}
    \Lambda_T
    \coloneqq
    \max_{1\le t\le T+1}\|\lambda_t\|_\infty,
    \qquad
    H\coloneqq 2L+E.
    \label{eq:app_constants}
\end{equation}
Then Assumption~\ref{ass:app_bounded} implies
\begin{equation}
    \|s_t\|_\infty\le L(1+\Lambda_T),
    \qquad
    \|h_t(p_t)\|_\infty\le H.
    \label{eq:app_basic_bounds}
\end{equation}
The second inequality follows from $\|u-p_t\|_1\le 2$ and $|[G_t]_{kj}|\le L$.

\subsection{Constrained no-regret guarantee}
\label{app:theory_regret}

For a fixed comparator $p\in\Delta_K$, define cumulative utility regret
\begin{equation}
    R_T(p)
    \coloneqq
    \sum_{t=1}^T\bigl(F_t(p)-F_t(p_t)\bigr).
    \label{eq:app_regret}
\end{equation}
The next theorem gives the finite-horizon bound from which the $O(\sqrt{T})$ statement in the main text follows.

\begin{theorem}[Constrained no-regret of full-Jacobian PDCM]
\label{thm:app_no_regret}
Suppose Assumptions~\ref{ass:app_concavity}--\ref{ass:app_bounded} hold and $(p_t,\lambda_t)$ is generated by Equations~\eqref{eq:app_primal_update}--\eqref{eq:app_dual_update}. Then for every horizon $T\ge 1$, every $p\in\mathcal{P}_T$, and every $\eta_p,\eta_\lambda>0$,
\begin{equation}
    R_T(p)
    \le
    \frac{\log K}{\eta_p}
    +
    \frac{\eta_p T}{2}L^2(1+\Lambda_T)^2
    +
    \frac{\eta_\lambda T}{2}KH^2.
    \label{eq:app_finite_regret_bound}
\end{equation}
Moreover, for every task $k$,
\begin{equation}
    \left[
        \sum_{t=1}^T c_{t,k}(p_t)
    \right]_+
    \le
    \left[
        \sum_{t=1}^T h_{t,k}(p_t)
    \right]_+
    \le
    \frac{\Lambda_T}{\eta_\lambda}.
    \label{eq:app_constraint_bound}
\end{equation}
If the dual sequence is stable, i.e., $\Lambda_T\le\Lambda<\infty$ for all $T$, and $\eta_p,\eta_\lambda=\Theta(T^{-1/2})$, then
\begin{equation}
    R_T(p)=O(\sqrt{T}),
    \qquad
    \left[
        \sum_{t=1}^T c_{t,k}(p_t)
    \right]_+
    =O(\sqrt{T}),
    \label{eq:app_sqrt_rates}
\end{equation}
so both average regret and average excess counterfactual lag vanish as $O(T^{-1/2})$.
\end{theorem}

\begin{proof}
We first prove a one-step entropic mirror-ascent inequality. By optimality of Equation~\eqref{eq:app_primal_update} and the three-point identity for the negative-entropy Bregman divergence, for every $p\in\Delta_K$,
\begin{equation}
    \eta_p\langle s_t,p-p_{t+1}\rangle
    \le
    D_{\mathrm{KL}}(p\|p_t)
    -D_{\mathrm{KL}}(p\|p_{t+1})
    -D_{\mathrm{KL}}(p_{t+1}\|p_t).
    \label{eq:app_three_point}
\end{equation}
Adding $\eta_p\langle s_t,p_{t+1}-p_t\rangle$ to both sides and using H\"older's inequality together with Pinsker's inequality,
\begin{equation}
    \langle s_t,p_{t+1}-p_t\rangle
    \le
    \|s_t\|_\infty\|p_{t+1}-p_t\|_1,
    \qquad
    D_{\mathrm{KL}}(p_{t+1}\|p_t)
    \ge
    \frac{1}{2}\|p_{t+1}-p_t\|_1^2,
\end{equation}
we obtain
\begin{equation}
    \langle s_t,p-p_t\rangle
    \le
    \frac{
        D_{\mathrm{KL}}(p\|p_t)
        -D_{\mathrm{KL}}(p\|p_{t+1})
    }{\eta_p}
    +
    \frac{\eta_p}{2}\|s_t\|_\infty^2.
    \label{eq:app_mirror_one_step}
\end{equation}
Summing over $t$, using $p_1=u$, $D_{\mathrm{KL}}(p\|u)\le\log K$, and Equation~\eqref{eq:app_basic_bounds}, gives
\begin{equation}
    \sum_{t=1}^T\langle s_t,p-p_t\rangle
    \le
    \frac{\log K}{\eta_p}
    +
    \frac{\eta_p T}{2}L^2(1+\Lambda_T)^2.
    \label{eq:app_mirror_sum}
\end{equation}

We now relate the score to utility regret and the safety residuals. From Equation~\eqref{eq:app_affine_residual},
\begin{equation}
    h_t(p_t)-h_t(p)=G_t(p-p_t).
    \label{eq:app_h_difference}
\end{equation}
Therefore, using Equation~\eqref{eq:app_score},
\begin{align}
    \langle s_t,p-p_t\rangle
    &=
    \left\langle G_t^\top\mathbf{1},p-p_t\right\rangle
    +
    \left\langle \lambda_t,G_t(p-p_t)\right\rangle\\
    &=
    \left\langle G_t^\top\mathbf{1},p-p_t\right\rangle
    +
    \left\langle \lambda_t,h_t(p_t)-h_t(p)\right\rangle.
    \label{eq:app_score_decomposition}
\end{align}
For $p\in\mathcal{P}_T$, we have $h_t(p)\le0$ coordinatewise and $\lambda_t\ge0$, so
\begin{equation}
    \left\langle G_t^\top\mathbf{1},p-p_t\right\rangle
    +
    \langle\lambda_t,h_t(p_t)\rangle
    \le
    \langle s_t,p-p_t\rangle.
    \label{eq:app_safe_drop}
\end{equation}
Since each $U_{t,k}$ is concave, $F_t$ is concave, and Lemma~\ref{lem:app_tangent} implies
\begin{equation}
    F_t(p)-F_t(p_t)
    \le
    \langle\nabla F_t(p_t),p-p_t\rangle
    =
    \left\langle G_t^\top\mathbf{1},p-p_t\right\rangle.
    \label{eq:app_F_concavity}
\end{equation}
Combining Equations~\eqref{eq:app_mirror_sum}, \eqref{eq:app_safe_drop}, and \eqref{eq:app_F_concavity} yields
\begin{equation}
    R_T(p)
    +
    \sum_{t=1}^T\langle\lambda_t,h_t(p_t)\rangle
    \le
    \frac{\log K}{\eta_p}
    +
    \frac{\eta_p T}{2}L^2(1+\Lambda_T)^2.
    \label{eq:app_predual}
\end{equation}

It remains to control the dual term. Projection onto the nonnegative orthant is nonexpansive and fixes the origin, hence Equation~\eqref{eq:app_dual_update} gives
\begin{align}
    \|\lambda_{t+1}\|_2^2
    &\le
    \|\lambda_t+\eta_\lambda h_t(p_t)\|_2^2\\
    &=
    \|\lambda_t\|_2^2
    +2\eta_\lambda\langle\lambda_t,h_t(p_t)\rangle
    +\eta_\lambda^2\|h_t(p_t)\|_2^2.
    \label{eq:app_dual_square}
\end{align}
Rearranging and summing from $t=1$ to $T$, using $\lambda_1=0$ and $\|h_t(p_t)\|_2^2\le KH^2$, gives
\begin{equation}
    -\sum_{t=1}^T\langle\lambda_t,h_t(p_t)\rangle
    \le
    \frac{\eta_\lambda T}{2}KH^2.
    \label{eq:app_dual_inner_bound}
\end{equation}
Substituting Equation~\eqref{eq:app_dual_inner_bound} into Equation~\eqref{eq:app_predual} proves Equation~\eqref{eq:app_finite_regret_bound}.

For the constraint bound, the coordinatewise projection satisfies
\begin{equation}
    \lambda_{t+1,k}
    =
    \max\{0,\lambda_{t,k}+\eta_\lambda h_{t,k}(p_t)\}
    \ge
    \lambda_{t,k}+\eta_\lambda h_{t,k}(p_t).
\end{equation}
Telescoping and using $\lambda_{1,k}=0$ gives
\begin{equation}
    \eta_\lambda\sum_{t=1}^T h_{t,k}(p_t)
    \le
    \lambda_{T+1,k}
    \le
    \Lambda_T.
\end{equation}
Taking positive parts proves the second inequality in Equation~\eqref{eq:app_constraint_bound}. The first follows from Equation~\eqref{eq:app_true_vs_linearized} after summing over $t$. Finally, if $\Lambda_T\le\Lambda$ and both step sizes are of order $T^{-1/2}$, every term in Equations~\eqref{eq:app_finite_regret_bound} and \eqref{eq:app_constraint_bound} is $O(\sqrt{T})$.
\end{proof}

The preceding theorem immediately yields the comparator used in the main text.

\begin{corollary}[Best fixed feasible allocation]
\label{cor:app_best_fixed}
Let
\begin{equation}
    p_T^\star
    \in
    \arg\max_{p\in\mathcal{P}_T}
    \sum_{t=1}^T F_t(p).
    \label{eq:app_best_fixed}
\end{equation}
Under the dual-stability condition $\Lambda_T\le\Lambda$ and step sizes $\eta_p,\eta_\lambda=\Theta(T^{-1/2})$,
\begin{equation}
    \sum_{t=1}^T
    \left[F_t(p_T^\star)-F_t(p_t)\right]
    =O(\sqrt{T}),
    \label{eq:app_best_fixed_regret}
\end{equation}
and for every task $k$,
\begin{equation}
    \left[
        \sum_{t=1}^T
        \left(
            U_{t,k}(u)-U_{t,k}(p_t)-\epsilon_{t,k}
        \right)
    \right]_+
    =O(\sqrt{T}).
    \label{eq:app_best_fixed_constraint}
\end{equation}
\end{corollary}

\begin{corollary}[Uniform-mixture safeguard]
\label{cor:app_uniform}
Because $u\in\mathcal{P}_T$, Theorem~\ref{thm:app_no_regret} applies in particular to the uniform mixture. Under dual stability and $\eta_p,\eta_\lambda=\Theta(T^{-1/2})$,
\begin{equation}
    \sum_{t=1}^T F_t(u)
    -
    \sum_{t=1}^T F_t(p_t)
    =O(\sqrt{T}).
    \label{eq:app_uniform_regret}
\end{equation}
Hence the average aggregate-utility gap between PDCM and the static uniform mixture vanishes as $O(T^{-1/2})$. In addition,
\begin{equation}
    \frac{1}{T}\sum_{t=1}^T
    \left(U_{t,k}(u)-U_{t,k}(p_t)\right)
    \le
    \frac{1}{T}\sum_{t=1}^T\epsilon_{t,k}
    +O(T^{-1/2})
    \label{eq:app_uniform_lag}
\end{equation}
for every task $k$.
\end{corollary}

\begin{remark}[Dual stability]
\label{rem:app_dual_stability}
The finite-horizon bounds in Equations~\eqref{eq:app_finite_regret_bound}--\eqref{eq:app_constraint_bound} hold for the realized quantity $\Lambda_T$ and do not assume that the dual sequence is bounded in advance. The $O(\sqrt{T})$ rates require $\Lambda_T=O(1)$. This is a stability condition on the coupled allocation and training dynamics; it may be checked empirically, imposed as an assumption, or enforced by an upper projection on the dual variables when such a projection is inactive relative to the desired optimum.
\end{remark}

\subsection{Dual-weighted marginal-utility matching}
\label{app:theory_matching}

The no-regret result describes finite-horizon behavior. A complementary stationary characterization explains the allocation condition targeted by the primal-dual coupling. Let $U_k:\Delta_K\to\mathbb{R}$ be differentiable and concave, define
\begin{equation}
    U(p)
    \coloneqq
    \begin{bmatrix}
        U_1(p)&\cdots&U_K(p)
    \end{bmatrix}^{\!\top},
    \qquad
    G(p)\coloneqq\nabla_p U(p),
    \qquad
    F(p)\coloneqq\mathbf{1}^\top U(p),
\end{equation}
and consider the static constrained allocation problem
\begin{equation}
\begin{aligned}
    \max_{p\in\Delta_K}\quad
    &F(p)\\
    \text{subject to}\quad
    &U_k(u)-U_k(p)-\epsilon_k\le0,
    \qquad k\in[K].
\end{aligned}
\label{eq:app_stationary_problem}
\end{equation}
Since $F$ is concave and each constraint function $U_k(u)-U_k(p)-\epsilon_k$ is convex, Equation~\eqref{eq:app_stationary_problem} is a convex optimization problem after negating the objective.

\begin{theorem}[Dual-weighted marginal-utility matching]
\label{thm:app_matching}
Assume Equation~\eqref{eq:app_stationary_problem} satisfies Slater's condition, and let $(p^\star,\lambda^\star)$ be a primal-dual optimum. Define
\begin{equation}
    s^\star
    \coloneqq
    G(p^\star)^\top(\mathbf{1}+\lambda^\star).
    \label{eq:app_stationary_score}
\end{equation}
Then there exists $\mu^\star\in\mathbb{R}$ such that
\begin{equation}
    s_j^\star=\mu^\star
    \quad\text{for every }j\text{ with }p_j^\star>0,
    \qquad
    s_j^\star\le\mu^\star
    \quad\text{for every }j\text{ with }p_j^\star=0.
    \label{eq:app_matching_condition}
\end{equation}
The multipliers additionally satisfy
\begin{equation}
    \lambda_k^\star\ge0,
    \qquad
    U_k(u)-U_k(p^\star)-\epsilon_k\le0,
    \qquad
    \lambda_k^\star\bigl(U_k(u)-U_k(p^\star)-\epsilon_k\bigr)=0.
    \label{eq:app_complementarity}
\end{equation}
If $p_j^\star>0$ for every $j$, then $p^\star$ is a fixed point of the exponentiated-gradient primal map with $\lambda^\star$ if and only if $s^\star=\mu^\star\mathbf{1}$.
\end{theorem}

\begin{proof}
Rewrite Equation~\eqref{eq:app_stationary_problem} as the convex minimization problem
\begin{equation}
\begin{aligned}
    \min_{p\in\Delta_K}\quad
    &-F(p)\\
    \text{subject to}\quad
    &c_k(p)\coloneqq U_k(u)-U_k(p)-\epsilon_k\le0,
    \qquad k\in[K].
\end{aligned}
\end{equation}
Under Slater's condition, the Karush--Kuhn--Tucker conditions are necessary and sufficient. Introduce multipliers $\lambda_k\ge0$ for $c_k(p)\le0$, $\mu\in\mathbb{R}$ for $\mathbf{1}^\top p=1$, and $\alpha_j\ge0$ for $-p_j\le0$. The KKT Lagrangian is
\begin{equation}
    \mathcal{L}_{\mathrm{KKT}}(p,\lambda,\mu,\alpha)
    =
    -F(p)
    +\sum_{k=1}^K\lambda_k\bigl(U_k(u)-U_k(p)-\epsilon_k\bigr)
    +\mu(\mathbf{1}^\top p-1)
    -\alpha^\top p.
\end{equation}
Differentiating with respect to $p$ gives the stationarity condition
\begin{equation}
    -G(p^\star)^\top(\mathbf{1}+\lambda^\star)
    +\mu^\star\mathbf{1}
    -\alpha^\star
    =0.
    \label{eq:app_KKT_stationarity}
\end{equation}
Hence $s_j^\star=\mu^\star-\alpha_j^\star$. If $p_j^\star>0$, complementary slackness for $-p_j\le0$ gives $\alpha_j^\star=0$, so $s_j^\star=\mu^\star$. If $p_j^\star=0$, then $\alpha_j^\star\ge0$, so $s_j^\star\le\mu^\star$. The remaining KKT conditions give Equation~\eqref{eq:app_complementarity}.

For the fixed-point statement, suppose first that $p_j^\star>0$ for every $j$ and $s^\star=\mu^\star\mathbf{1}$. Then the exponentiated-gradient map gives
\begin{equation}
    \frac{
        p_j^\star\exp(\eta_p s_j^\star)
    }{
        \sum_{\ell=1}^K p_\ell^\star\exp(\eta_p s_\ell^\star)
    }
    =
    \frac{
        p_j^\star\exp(\eta_p\mu^\star)
    }{
        \exp(\eta_p\mu^\star)\sum_{\ell=1}^K p_\ell^\star
    }
    =p_j^\star.
\end{equation}
Conversely, if the map fixes an interior $p^\star$, division by $p_j^\star>0$ shows that $\exp(\eta_p s_j^\star)$ equals the same normalization constant for every $j$. Taking logarithms gives $s_j^\star=s_\ell^\star$ for all $j,\ell$, which is exactly $s^\star=\mu^\star\mathbf{1}$ for some scalar $\mu^\star$.
\end{proof}

Theorem~\ref{thm:app_matching} exposes the role of the full Jacobian. Expanding the stationary score coordinatewise,
\begin{equation}
    s_j^\star
    =
    \sum_{k=1}^K
    (1+\lambda_k^\star)
    \frac{\partial U_k(p^\star)}{\partial p_j}.
    \label{eq:app_cross_task_matching}
\end{equation}
Thus the dual variable of task $k$ amplifies the entire $k$th row of the allocation-response Jacobian, not only its diagonal entry. A training coordinate $j$ is valuable when increasing its allocation improves any evaluation task with high marginal utility, and its score receives additional weight when those tasks are under active lag pressure. At an interior stationary optimum, these dual-weighted cross-task marginal effects are equalized across all training coordinates.
 from the main document.
% Requires in the preamble of the including file:
%   \usepackage{amsthm,mathtools}
%   \newtheorem{assumption}{Assumption}
%   \newtheorem{lemma}{Lemma}
%   \newtheorem{theorem}{Theorem}
%   \newtheorem{corollary}{Corollary}
%   \newtheorem{remark}{Remark}
%   \newcommand{\DeltaK}{\Delta_K}
%   \newcommand{\one}{\mathbf{1}}
%   \newcommand{\pos}[1]{\left[#1\right]_+}
%   \newcommand{\ip}[2]{\left\langle #1,#2\right\rangle}
%   \newcommand{\norm}[1]{\left\lVert #1\right\rVert}
% \R and \KL come from math_commands.tex (lines 480, 488).

\section{Theory}
\label{app:theory}

This appendix gives a self-contained analysis of the full-Jacobian version of PDCM used in Algorithm~\ref{alg:pdcm_jacobian}. The analysis isolates the online allocation mechanism. At each round it treats the allocation-response Jacobian as available exactly.

\subsection{Problem formulation and idealized full-Jacobian PDCM}
\label{app:theory_setup}

We consider $K\ge 2$ tasks and the probability simplex
\begin{equation}
    \Delta_K
    \coloneqq
    \left\{p\in\mathbb{R}_+^K:\mathbf{1}^\top p=1\right\},
    \qquad
    u\coloneqq \frac{1}{K}\mathbf{1}.
\end{equation}
At round $t$, let
\begin{equation}
    U_t(p)
    \coloneqq
    \begin{bmatrix}
        U_{t,1}(p) & \cdots & U_{t,K}(p)
    \end{bmatrix}^{\!\top}\in\mathbb{R}^K
\end{equation}
denote the vector of expected one-step policy improvements that would result from applying the round-$t$ update under allocation $p\in\Delta_K$. The aggregate improvement is
\begin{equation}
    F_t(p)\coloneqq \mathbf{1}^\top U_t(p)=\sum_{k=1}^K U_{t,k}(p).
    \label{eq:app_aggregate_utility}
\end{equation}
Because the policy state changes with training, the functions $U_t$ may be nonstationary in $t$.

We use the following structural condition. It is stronger than the reference-direction inequality needed only for the safety certificate, but it is the condition required to compare PDCM with an arbitrary fixed allocation in the regret statement.

\begin{assumption}[Concave allocation response]
\label{ass:app_concavity}
For every round $t$ and task $k$, $U_{t,k}$ is differentiable and concave on $\Delta_K$. Equivalently, for every $p,q\in\Delta_K$ and $\alpha\in[0,1]$,
\begin{equation}
    U_{t,k}((1-\alpha)p+\alpha q)
    \ge
    (1-\alpha)U_{t,k}(p)+\alpha U_{t,k}(q).
    \label{eq:app_line_concavity}
\end{equation}
\end{assumption}

Assumption~\ref{ass:app_concavity} implies the usual first-order upper bound for a concave function.

\begin{lemma}[Tangent upper bound]
\label{lem:app_tangent}
Under Assumption~\ref{ass:app_concavity}, for every $t,k$ and $p,q\in\Delta_K$,
\begin{equation}
    U_{t,k}(q)-U_{t,k}(p)
    \le
    \left\langle \nabla U_{t,k}(p),q-p\right\rangle.
    \label{eq:app_tangent_bound}
\end{equation}
\end{lemma}

\begin{proof}
For $\alpha\in(0,1]$, concavity gives
\begin{equation}
    U_{t,k}\!\left(p+\alpha(q-p)\right)-U_{t,k}(p)
    \ge
    \alpha\bigl(U_{t,k}(q)-U_{t,k}(p)\bigr).
\end{equation}
Divide by $\alpha$ and let $\alpha\downarrow 0$. Differentiability yields Equation~\eqref{eq:app_tangent_bound}.
\end{proof}

At the played allocation $p_t$, define the allocation-response Jacobian
\begin{equation}
    G_t
    \coloneqq
    \nabla_p U_t(p_t)
    =
    \begin{bmatrix}
        g_{t,1}^{\top}\\
        \vdots\\
        g_{t,K}^{\top}
    \end{bmatrix}
    \in\mathbb{R}^{K\times K},
    \qquad
    g_{t,k}\coloneqq \nabla_p U_{t,k}(p_t).
    \label{eq:app_jacobian}
\end{equation}
Thus $[G_t]_{kj}=\partial U_{t,k}(p_t)/\partial p_j$ measures the local effect of allocating additional mass to training task $j$ on the improvement of evaluation task $k$. In particular,
\begin{equation}
    \nabla F_t(p_t)=G_t^\top\mathbf{1}.
    \label{eq:app_aggregate_gradient}
\end{equation}

For task $k$, let $\epsilon_{t,k}\ge 0$ denote the lag budget assigned to round $t$. For a total cumulative budget $\epsilon_k$, one may choose any schedule satisfying $\sum_{t=1}^T\epsilon_{t,k}=\epsilon_k$, for example $\epsilon_{t,k}=\epsilon_k/T$. Define the true one-step counterfactual lag residual
\begin{equation}
    c_{t,k}(p_t)
    \coloneqq
    U_{t,k}(u)-U_{t,k}(p_t)-\epsilon_{t,k}.
    \label{eq:app_true_residual}
\end{equation}
Applying Lemma~\ref{lem:app_tangent} with $p=p_t$ and $q=u$ gives
\begin{equation}
    c_{t,k}(p_t)
    \le
    \left\langle g_{t,k},u-p_t\right\rangle-\epsilon_{t,k}.
    \label{eq:app_safety_certificate}
\end{equation}
Therefore the first-order quantity on the right is a conservative certificate of the unobserved uniform-allocation lag.

For the online analysis, extend this certificate affinely to an arbitrary comparator $p\in\Delta_K$:
\begin{equation}
    h_t(p)
    \coloneqq
    G_t(u-p)-\epsilon_t,
    \qquad
    \epsilon_t
    \coloneqq
    (\epsilon_{t,1},\ldots,\epsilon_{t,K})^\top.
    \label{eq:app_affine_residual}
\end{equation}
Equation~\eqref{eq:app_safety_certificate} is equivalently
\begin{equation}
    c_t(p_t)\le h_t(p_t)
    \qquad\text{coordinatewise}.
    \label{eq:app_true_vs_linearized}
\end{equation}
For horizon $T$, define the fixed first-order-safe comparator set
\begin{equation}
    \mathcal{P}_T
    \coloneqq
    \left\{
        p\in\Delta_K:
        h_t(p)\le 0\ \text{coordinatewise for every }t\le T
    \right\}.
    \label{eq:app_comparator_set}
\end{equation}
This set is always nonempty because $h_t(u)=-\epsilon_t\le 0$.

Let $\lambda_t\in\mathbb{R}_+^K$ be the task-wise dual variables. Dropping terms independent of $p$ from the linearized round-$t$ Lagrangian gives the primal score
\begin{equation}
    s_t
    \coloneqq
    G_t^\top(\mathbf{1}+\lambda_t).
    \label{eq:app_score}
\end{equation}
The idealized full-Jacobian PDCM recursion is
\begin{align}
    p_{t+1}
    &\coloneqq
    \arg\max_{p\in\Delta_K}
    \left\{
        \eta_p\langle s_t,p\rangle
        -D_{\mathrm{KL}}(p\|p_t)
    \right\},
    \label{eq:app_primal_update}\\
    \lambda_{t+1}
    &\coloneqq
    \left[
        \lambda_t+\eta_\lambda h_t(p_t)
    \right]_+,
    \label{eq:app_dual_update}
\end{align}
with $p_1=u$ and $\lambda_1=0$. Here $[\cdot]_+$ denotes coordinatewise projection onto $\mathbb{R}_+^K$. The primal step has the closed form
\begin{equation}
    p_{t+1,j}
    =
    \frac{
        p_{t,j}\exp(\eta_p s_{t,j})
    }{
        \sum_{\ell=1}^K p_{t,\ell}\exp(\eta_p s_{t,\ell})
    }.
    \label{eq:app_exp_update}
\end{equation}

We impose only boundedness conditions needed by the finite-horizon analysis.

\begin{assumption}[Bounded Jacobian and lag budget]
\label{ass:app_bounded}
There exist finite constants $L,E$ such that for every round $t$,
\begin{equation}
    \|G_t\|_1
    \coloneqq
    \max_{j\in[K]}\sum_{k=1}^K |[G_t]_{kj}|
    \le L,
    \qquad
    0\le \epsilon_{t,k}\le E.
    \label{eq:app_boundedness}
\end{equation}
\end{assumption}

For a horizon $T$, define the realized dual magnitude
\begin{equation}
    \Lambda_T
    \coloneqq
    \max_{1\le t\le T+1}\|\lambda_t\|_\infty,
    \qquad
    H\coloneqq 2L+E.
    \label{eq:app_constants}
\end{equation}
Then Assumption~\ref{ass:app_bounded} implies
\begin{equation}
    \|s_t\|_\infty\le L(1+\Lambda_T),
    \qquad
    \|h_t(p_t)\|_\infty\le H.
    \label{eq:app_basic_bounds}
\end{equation}
The second inequality follows from $\|u-p_t\|_1\le 2$ and $|[G_t]_{kj}|\le L$.

\subsection{Constrained no-regret guarantee}
\label{app:theory_regret}

For a fixed comparator $p\in\Delta_K$, define cumulative utility regret
\begin{equation}
    R_T(p)
    \coloneqq
    \sum_{t=1}^T\bigl(F_t(p)-F_t(p_t)\bigr).
    \label{eq:app_regret}
\end{equation}
The next theorem gives the finite-horizon bound from which the $O(\sqrt{T})$ statement in the main text follows.

\begin{theorem}[Constrained no-regret of full-Jacobian PDCM]
\label{thm:app_no_regret}
Suppose Assumptions~\ref{ass:app_concavity}--\ref{ass:app_bounded} hold and $(p_t,\lambda_t)$ is generated by Equations~\eqref{eq:app_primal_update}--\eqref{eq:app_dual_update}. Then for every horizon $T\ge 1$, every $p\in\mathcal{P}_T$, and every $\eta_p,\eta_\lambda>0$,
\begin{equation}
    R_T(p)
    \le
    \frac{\log K}{\eta_p}
    +
    \frac{\eta_p T}{2}L^2(1+\Lambda_T)^2
    +
    \frac{\eta_\lambda T}{2}KH^2.
    \label{eq:app_finite_regret_bound}
\end{equation}
Moreover, for every task $k$,
\begin{equation}
    \left[
        \sum_{t=1}^T c_{t,k}(p_t)
    \right]_+
    \le
    \left[
        \sum_{t=1}^T h_{t,k}(p_t)
    \right]_+
    \le
    \frac{\Lambda_T}{\eta_\lambda}.
    \label{eq:app_constraint_bound}
\end{equation}
If the dual sequence is stable, i.e., $\Lambda_T\le\Lambda<\infty$ for all $T$, and $\eta_p,\eta_\lambda=\Theta(T^{-1/2})$, then
\begin{equation}
    R_T(p)=O(\sqrt{T}),
    \qquad
    \left[
        \sum_{t=1}^T c_{t,k}(p_t)
    \right]_+
    =O(\sqrt{T}),
    \label{eq:app_sqrt_rates}
\end{equation}
so both average regret and average excess counterfactual lag vanish as $O(T^{-1/2})$.
\end{theorem}

\begin{proof}
We first prove a one-step entropic mirror-ascent inequality. By optimality of Equation~\eqref{eq:app_primal_update} and the three-point identity for the negative-entropy Bregman divergence, for every $p\in\Delta_K$,
\begin{equation}
    \eta_p\langle s_t,p-p_{t+1}\rangle
    \le
    D_{\mathrm{KL}}(p\|p_t)
    -D_{\mathrm{KL}}(p\|p_{t+1})
    -D_{\mathrm{KL}}(p_{t+1}\|p_t).
    \label{eq:app_three_point}
\end{equation}
Adding $\eta_p\langle s_t,p_{t+1}-p_t\rangle$ to both sides and using H\"older's inequality together with Pinsker's inequality,
\begin{equation}
    \langle s_t,p_{t+1}-p_t\rangle
    \le
    \|s_t\|_\infty\|p_{t+1}-p_t\|_1,
    \qquad
    D_{\mathrm{KL}}(p_{t+1}\|p_t)
    \ge
    \frac{1}{2}\|p_{t+1}-p_t\|_1^2,
\end{equation}
we obtain
\begin{equation}
    \langle s_t,p-p_t\rangle
    \le
    \frac{
        D_{\mathrm{KL}}(p\|p_t)
        -D_{\mathrm{KL}}(p\|p_{t+1})
    }{\eta_p}
    +
    \frac{\eta_p}{2}\|s_t\|_\infty^2.
    \label{eq:app_mirror_one_step}
\end{equation}
Summing over $t$, using $p_1=u$, $D_{\mathrm{KL}}(p\|u)\le\log K$, and Equation~\eqref{eq:app_basic_bounds}, gives
\begin{equation}
    \sum_{t=1}^T\langle s_t,p-p_t\rangle
    \le
    \frac{\log K}{\eta_p}
    +
    \frac{\eta_p T}{2}L^2(1+\Lambda_T)^2.
    \label{eq:app_mirror_sum}
\end{equation}

We now relate the score to utility regret and the safety residuals. From Equation~\eqref{eq:app_affine_residual},
\begin{equation}
    h_t(p_t)-h_t(p)=G_t(p-p_t).
    \label{eq:app_h_difference}
\end{equation}
Therefore, using Equation~\eqref{eq:app_score},
\begin{align}
    \langle s_t,p-p_t\rangle
    &=
    \left\langle G_t^\top\mathbf{1},p-p_t\right\rangle
    +
    \left\langle \lambda_t,G_t(p-p_t)\right\rangle\\
    &=
    \left\langle G_t^\top\mathbf{1},p-p_t\right\rangle
    +
    \left\langle \lambda_t,h_t(p_t)-h_t(p)\right\rangle.
    \label{eq:app_score_decomposition}
\end{align}
For $p\in\mathcal{P}_T$, we have $h_t(p)\le0$ coordinatewise and $\lambda_t\ge0$, so
\begin{equation}
    \left\langle G_t^\top\mathbf{1},p-p_t\right\rangle
    +
    \langle\lambda_t,h_t(p_t)\rangle
    \le
    \langle s_t,p-p_t\rangle.
    \label{eq:app_safe_drop}
\end{equation}
Since each $U_{t,k}$ is concave, $F_t$ is concave, and Lemma~\ref{lem:app_tangent} implies
\begin{equation}
    F_t(p)-F_t(p_t)
    \le
    \langle\nabla F_t(p_t),p-p_t\rangle
    =
    \left\langle G_t^\top\mathbf{1},p-p_t\right\rangle.
    \label{eq:app_F_concavity}
\end{equation}
Combining Equations~\eqref{eq:app_mirror_sum}, \eqref{eq:app_safe_drop}, and \eqref{eq:app_F_concavity} yields
\begin{equation}
    R_T(p)
    +
    \sum_{t=1}^T\langle\lambda_t,h_t(p_t)\rangle
    \le
    \frac{\log K}{\eta_p}
    +
    \frac{\eta_p T}{2}L^2(1+\Lambda_T)^2.
    \label{eq:app_predual}
\end{equation}

It remains to control the dual term. Projection onto the nonnegative orthant is nonexpansive and fixes the origin, hence Equation~\eqref{eq:app_dual_update} gives
\begin{align}
    \|\lambda_{t+1}\|_2^2
    &\le
    \|\lambda_t+\eta_\lambda h_t(p_t)\|_2^2\\
    &=
    \|\lambda_t\|_2^2
    +2\eta_\lambda\langle\lambda_t,h_t(p_t)\rangle
    +\eta_\lambda^2\|h_t(p_t)\|_2^2.
    \label{eq:app_dual_square}
\end{align}
Rearranging and summing from $t=1$ to $T$, using $\lambda_1=0$ and $\|h_t(p_t)\|_2^2\le KH^2$, gives
\begin{equation}
    -\sum_{t=1}^T\langle\lambda_t,h_t(p_t)\rangle
    \le
    \frac{\eta_\lambda T}{2}KH^2.
    \label{eq:app_dual_inner_bound}
\end{equation}
Substituting Equation~\eqref{eq:app_dual_inner_bound} into Equation~\eqref{eq:app_predual} proves Equation~\eqref{eq:app_finite_regret_bound}.

For the constraint bound, the coordinatewise projection satisfies
\begin{equation}
    \lambda_{t+1,k}
    =
    \max\{0,\lambda_{t,k}+\eta_\lambda h_{t,k}(p_t)\}
    \ge
    \lambda_{t,k}+\eta_\lambda h_{t,k}(p_t).
\end{equation}
Telescoping and using $\lambda_{1,k}=0$ gives
\begin{equation}
    \eta_\lambda\sum_{t=1}^T h_{t,k}(p_t)
    \le
    \lambda_{T+1,k}
    \le
    \Lambda_T.
\end{equation}
Taking positive parts proves the second inequality in Equation~\eqref{eq:app_constraint_bound}. The first follows from Equation~\eqref{eq:app_true_vs_linearized} after summing over $t$. Finally, if $\Lambda_T\le\Lambda$ and both step sizes are of order $T^{-1/2}$, every term in Equations~\eqref{eq:app_finite_regret_bound} and \eqref{eq:app_constraint_bound} is $O(\sqrt{T})$.
\end{proof}

The preceding theorem immediately yields the comparator used in the main text.

\begin{corollary}[Best fixed feasible allocation]
\label{cor:app_best_fixed}
Let
\begin{equation}
    p_T^\star
    \in
    \arg\max_{p\in\mathcal{P}_T}
    \sum_{t=1}^T F_t(p).
    \label{eq:app_best_fixed}
\end{equation}
Under the dual-stability condition $\Lambda_T\le\Lambda$ and step sizes $\eta_p,\eta_\lambda=\Theta(T^{-1/2})$,
\begin{equation}
    \sum_{t=1}^T
    \left[F_t(p_T^\star)-F_t(p_t)\right]
    =O(\sqrt{T}),
    \label{eq:app_best_fixed_regret}
\end{equation}
and for every task $k$,
\begin{equation}
    \left[
        \sum_{t=1}^T
        \left(
            U_{t,k}(u)-U_{t,k}(p_t)-\epsilon_{t,k}
        \right)
    \right]_+
    =O(\sqrt{T}).
    \label{eq:app_best_fixed_constraint}
\end{equation}
\end{corollary}

\begin{corollary}[Uniform-mixture safeguard]
\label{cor:app_uniform}
Because $u\in\mathcal{P}_T$, Theorem~\ref{thm:app_no_regret} applies in particular to the uniform mixture. Under dual stability and $\eta_p,\eta_\lambda=\Theta(T^{-1/2})$,
\begin{equation}
    \sum_{t=1}^T F_t(u)
    -
    \sum_{t=1}^T F_t(p_t)
    =O(\sqrt{T}).
    \label{eq:app_uniform_regret}
\end{equation}
Hence the average aggregate-utility gap between PDCM and the static uniform mixture vanishes as $O(T^{-1/2})$. In addition,
\begin{equation}
    \frac{1}{T}\sum_{t=1}^T
    \left(U_{t,k}(u)-U_{t,k}(p_t)\right)
    \le
    \frac{1}{T}\sum_{t=1}^T\epsilon_{t,k}
    +O(T^{-1/2})
    \label{eq:app_uniform_lag}
\end{equation}
for every task $k$.
\end{corollary}

\begin{remark}[Dual stability]
\label{rem:app_dual_stability}
The finite-horizon bounds in Equations~\eqref{eq:app_finite_regret_bound}--\eqref{eq:app_constraint_bound} hold for the realized quantity $\Lambda_T$ and do not assume that the dual sequence is bounded in advance. The $O(\sqrt{T})$ rates require $\Lambda_T=O(1)$. This is a stability condition on the coupled allocation and training dynamics; it may be checked empirically, imposed as an assumption, or enforced by an upper projection on the dual variables when such a projection is inactive relative to the desired optimum.
\end{remark}

\subsection{Dual-weighted marginal-utility matching}
\label{app:theory_matching}

The no-regret result describes finite-horizon behavior. A complementary stationary characterization explains the allocation condition targeted by the primal-dual coupling. Let $U_k:\Delta_K\to\mathbb{R}$ be differentiable and concave, define
\begin{equation}
    U(p)
    \coloneqq
    \begin{bmatrix}
        U_1(p)&\cdots&U_K(p)
    \end{bmatrix}^{\!\top},
    \qquad
    G(p)\coloneqq\nabla_p U(p),
    \qquad
    F(p)\coloneqq\mathbf{1}^\top U(p),
\end{equation}
and consider the static constrained allocation problem
\begin{equation}
\begin{aligned}
    \max_{p\in\Delta_K}\quad
    &F(p)\\
    \text{subject to}\quad
    &U_k(u)-U_k(p)-\epsilon_k\le0,
    \qquad k\in[K].
\end{aligned}
\label{eq:app_stationary_problem}
\end{equation}
Since $F$ is concave and each constraint function $U_k(u)-U_k(p)-\epsilon_k$ is convex, Equation~\eqref{eq:app_stationary_problem} is a convex optimization problem after negating the objective.

\begin{theorem}[Dual-weighted marginal-utility matching]
\label{thm:app_matching}
Assume Equation~\eqref{eq:app_stationary_problem} satisfies Slater's condition, and let $(p^\star,\lambda^\star)$ be a primal-dual optimum. Define
\begin{equation}
    s^\star
    \coloneqq
    G(p^\star)^\top(\mathbf{1}+\lambda^\star).
    \label{eq:app_stationary_score}
\end{equation}
Then there exists $\mu^\star\in\mathbb{R}$ such that
\begin{equation}
    s_j^\star=\mu^\star
    \quad\text{for every }j\text{ with }p_j^\star>0,
    \qquad
    s_j^\star\le\mu^\star
    \quad\text{for every }j\text{ with }p_j^\star=0.
    \label{eq:app_matching_condition}
\end{equation}
The multipliers additionally satisfy
\begin{equation}
    \lambda_k^\star\ge0,
    \qquad
    U_k(u)-U_k(p^\star)-\epsilon_k\le0,
    \qquad
    \lambda_k^\star\bigl(U_k(u)-U_k(p^\star)-\epsilon_k\bigr)=0.
    \label{eq:app_complementarity}
\end{equation}
If $p_j^\star>0$ for every $j$, then $p^\star$ is a fixed point of the exponentiated-gradient primal map with $\lambda^\star$ if and only if $s^\star=\mu^\star\mathbf{1}$.
\end{theorem}

\begin{proof}
Rewrite Equation~\eqref{eq:app_stationary_problem} as the convex minimization problem
\begin{equation}
\begin{aligned}
    \min_{p\in\Delta_K}\quad
    &-F(p)\\
    \text{subject to}\quad
    &c_k(p)\coloneqq U_k(u)-U_k(p)-\epsilon_k\le0,
    \qquad k\in[K].
\end{aligned}
\end{equation}
Under Slater's condition, the Karush--Kuhn--Tucker conditions are necessary and sufficient. Introduce multipliers $\lambda_k\ge0$ for $c_k(p)\le0$, $\mu\in\mathbb{R}$ for $\mathbf{1}^\top p=1$, and $\alpha_j\ge0$ for $-p_j\le0$. The KKT Lagrangian is
\begin{equation}
    \mathcal{L}_{\mathrm{KKT}}(p,\lambda,\mu,\alpha)
    =
    -F(p)
    +\sum_{k=1}^K\lambda_k\bigl(U_k(u)-U_k(p)-\epsilon_k\bigr)
    +\mu(\mathbf{1}^\top p-1)
    -\alpha^\top p.
\end{equation}
Differentiating with respect to $p$ gives the stationarity condition
\begin{equation}
    -G(p^\star)^\top(\mathbf{1}+\lambda^\star)
    +\mu^\star\mathbf{1}
    -\alpha^\star
    =0.
    \label{eq:app_KKT_stationarity}
\end{equation}
Hence $s_j^\star=\mu^\star-\alpha_j^\star$. If $p_j^\star>0$, complementary slackness for $-p_j\le0$ gives $\alpha_j^\star=0$, so $s_j^\star=\mu^\star$. If $p_j^\star=0$, then $\alpha_j^\star\ge0$, so $s_j^\star\le\mu^\star$. The remaining KKT conditions give Equation~\eqref{eq:app_complementarity}.

For the fixed-point statement, suppose first that $p_j^\star>0$ for every $j$ and $s^\star=\mu^\star\mathbf{1}$. Then the exponentiated-gradient map gives
\begin{equation}
    \frac{
        p_j^\star\exp(\eta_p s_j^\star)
    }{
        \sum_{\ell=1}^K p_\ell^\star\exp(\eta_p s_\ell^\star)
    }
    =
    \frac{
        p_j^\star\exp(\eta_p\mu^\star)
    }{
        \exp(\eta_p\mu^\star)\sum_{\ell=1}^K p_\ell^\star
    }
    =p_j^\star.
\end{equation}
Conversely, if the map fixes an interior $p^\star$, division by $p_j^\star>0$ shows that $\exp(\eta_p s_j^\star)$ equals the same normalization constant for every $j$. Taking logarithms gives $s_j^\star=s_\ell^\star$ for all $j,\ell$, which is exactly $s^\star=\mu^\star\mathbf{1}$ for some scalar $\mu^\star$.
\end{proof}

Theorem~\ref{thm:app_matching} exposes the role of the full Jacobian. Expanding the stationary score coordinatewise,
\begin{equation}
    s_j^\star
    =
    \sum_{k=1}^K
    (1+\lambda_k^\star)
    \frac{\partial U_k(p^\star)}{\partial p_j}.
    \label{eq:app_cross_task_matching}
\end{equation}
Thus the dual variable of task $k$ amplifies the entire $k$th row of the allocation-response Jacobian, not only its diagonal entry. A training coordinate $j$ is valuable when increasing its allocation improves any evaluation task with high marginal utility, and its score receives additional weight when those tasks are under active lag pressure. At an interior stationary optimum, these dual-weighted cross-task marginal effects are equalized across all training coordinates.
 from the main document.
% Requires in the preamble of the including file:
%   \usepackage{amsthm,mathtools}
%   \newtheorem{assumption}{Assumption}
%   \newtheorem{lemma}{Lemma}
%   \newtheorem{theorem}{Theorem}
%   \newtheorem{corollary}{Corollary}
%   \newtheorem{remark}{Remark}
%   \newcommand{\DeltaK}{\Delta_K}
%   \newcommand{\one}{\mathbf{1}}
%   \newcommand{\pos}[1]{\left[#1\right]_+}
%   \newcommand{\ip}[2]{\left\langle #1,#2\right\rangle}
%   \newcommand{\norm}[1]{\left\lVert #1\right\rVert}
% \R and \KL come from math_commands.tex (lines 480, 488).

\section{Theory}
\label{app:theory}

This appendix gives a self-contained analysis of the full-Jacobian version of PDCM used in Algorithm~\ref{alg:pdcm_jacobian}. The analysis isolates the online allocation mechanism. At each round it treats the allocation-response Jacobian as available exactly.

\subsection{Problem formulation and idealized full-Jacobian PDCM}
\label{app:theory_setup}

We consider $K\ge 2$ tasks and the probability simplex
\begin{equation}
    \Delta_K
    \coloneqq
    \left\{p\in\mathbb{R}_+^K:\mathbf{1}^\top p=1\right\},
    \qquad
    u\coloneqq \frac{1}{K}\mathbf{1}.
\end{equation}
At round $t$, let
\begin{equation}
    U_t(p)
    \coloneqq
    \begin{bmatrix}
        U_{t,1}(p) & \cdots & U_{t,K}(p)
    \end{bmatrix}^{\!\top}\in\mathbb{R}^K
\end{equation}
denote the vector of expected one-step policy improvements that would result from applying the round-$t$ update under allocation $p\in\Delta_K$. The aggregate improvement is
\begin{equation}
    F_t(p)\coloneqq \mathbf{1}^\top U_t(p)=\sum_{k=1}^K U_{t,k}(p).
    \label{eq:app_aggregate_utility}
\end{equation}
Because the policy state changes with training, the functions $U_t$ may be nonstationary in $t$.

We use the following structural condition. It is stronger than the reference-direction inequality needed only for the safety certificate, but it is the condition required to compare PDCM with an arbitrary fixed allocation in the regret statement.

\begin{assumption}[Concave allocation response]
\label{ass:app_concavity}
For every round $t$ and task $k$, $U_{t,k}$ is differentiable and concave on $\Delta_K$. Equivalently, for every $p,q\in\Delta_K$ and $\alpha\in[0,1]$,
\begin{equation}
    U_{t,k}((1-\alpha)p+\alpha q)
    \ge
    (1-\alpha)U_{t,k}(p)+\alpha U_{t,k}(q).
    \label{eq:app_line_concavity}
\end{equation}
\end{assumption}

Assumption~\ref{ass:app_concavity} implies the usual first-order upper bound for a concave function.

\begin{lemma}[Tangent upper bound]
\label{lem:app_tangent}
Under Assumption~\ref{ass:app_concavity}, for every $t,k$ and $p,q\in\Delta_K$,
\begin{equation}
    U_{t,k}(q)-U_{t,k}(p)
    \le
    \left\langle \nabla U_{t,k}(p),q-p\right\rangle.
    \label{eq:app_tangent_bound}
\end{equation}
\end{lemma}

\begin{proof}
For $\alpha\in(0,1]$, concavity gives
\begin{equation}
    U_{t,k}\!\left(p+\alpha(q-p)\right)-U_{t,k}(p)
    \ge
    \alpha\bigl(U_{t,k}(q)-U_{t,k}(p)\bigr).
\end{equation}
Divide by $\alpha$ and let $\alpha\downarrow 0$. Differentiability yields Equation~\eqref{eq:app_tangent_bound}.
\end{proof}

At the played allocation $p_t$, define the allocation-response Jacobian
\begin{equation}
    G_t
    \coloneqq
    \nabla_p U_t(p_t)
    =
    \begin{bmatrix}
        g_{t,1}^{\top}\\
        \vdots\\
        g_{t,K}^{\top}
    \end{bmatrix}
    \in\mathbb{R}^{K\times K},
    \qquad
    g_{t,k}\coloneqq \nabla_p U_{t,k}(p_t).
    \label{eq:app_jacobian}
\end{equation}
Thus $[G_t]_{kj}=\partial U_{t,k}(p_t)/\partial p_j$ measures the local effect of allocating additional mass to training task $j$ on the improvement of evaluation task $k$. In particular,
\begin{equation}
    \nabla F_t(p_t)=G_t^\top\mathbf{1}.
    \label{eq:app_aggregate_gradient}
\end{equation}

For task $k$, let $\epsilon_{t,k}\ge 0$ denote the lag budget assigned to round $t$. For a total cumulative budget $\epsilon_k$, one may choose any schedule satisfying $\sum_{t=1}^T\epsilon_{t,k}=\epsilon_k$, for example $\epsilon_{t,k}=\epsilon_k/T$. Define the true one-step counterfactual lag residual
\begin{equation}
    c_{t,k}(p_t)
    \coloneqq
    U_{t,k}(u)-U_{t,k}(p_t)-\epsilon_{t,k}.
    \label{eq:app_true_residual}
\end{equation}
Applying Lemma~\ref{lem:app_tangent} with $p=p_t$ and $q=u$ gives
\begin{equation}
    c_{t,k}(p_t)
    \le
    \left\langle g_{t,k},u-p_t\right\rangle-\epsilon_{t,k}.
    \label{eq:app_safety_certificate}
\end{equation}
Therefore the first-order quantity on the right is a conservative certificate of the unobserved uniform-allocation lag.

For the online analysis, extend this certificate affinely to an arbitrary comparator $p\in\Delta_K$:
\begin{equation}
    h_t(p)
    \coloneqq
    G_t(u-p)-\epsilon_t,
    \qquad
    \epsilon_t
    \coloneqq
    (\epsilon_{t,1},\ldots,\epsilon_{t,K})^\top.
    \label{eq:app_affine_residual}
\end{equation}
Equation~\eqref{eq:app_safety_certificate} is equivalently
\begin{equation}
    c_t(p_t)\le h_t(p_t)
    \qquad\text{coordinatewise}.
    \label{eq:app_true_vs_linearized}
\end{equation}
For horizon $T$, define the fixed first-order-safe comparator set
\begin{equation}
    \mathcal{P}_T
    \coloneqq
    \left\{
        p\in\Delta_K:
        h_t(p)\le 0\ \text{coordinatewise for every }t\le T
    \right\}.
    \label{eq:app_comparator_set}
\end{equation}
This set is always nonempty because $h_t(u)=-\epsilon_t\le 0$.

Let $\lambda_t\in\mathbb{R}_+^K$ be the task-wise dual variables. Dropping terms independent of $p$ from the linearized round-$t$ Lagrangian gives the primal score
\begin{equation}
    s_t
    \coloneqq
    G_t^\top(\mathbf{1}+\lambda_t).
    \label{eq:app_score}
\end{equation}
The idealized full-Jacobian PDCM recursion is
\begin{align}
    p_{t+1}
    &\coloneqq
    \arg\max_{p\in\Delta_K}
    \left\{
        \eta_p\langle s_t,p\rangle
        -D_{\mathrm{KL}}(p\|p_t)
    \right\},
    \label{eq:app_primal_update}\\
    \lambda_{t+1}
    &\coloneqq
    \left[
        \lambda_t+\eta_\lambda h_t(p_t)
    \right]_+,
    \label{eq:app_dual_update}
\end{align}
with $p_1=u$ and $\lambda_1=0$. Here $[\cdot]_+$ denotes coordinatewise projection onto $\mathbb{R}_+^K$. The primal step has the closed form
\begin{equation}
    p_{t+1,j}
    =
    \frac{
        p_{t,j}\exp(\eta_p s_{t,j})
    }{
        \sum_{\ell=1}^K p_{t,\ell}\exp(\eta_p s_{t,\ell})
    }.
    \label{eq:app_exp_update}
\end{equation}

We impose only boundedness conditions needed by the finite-horizon analysis.

\begin{assumption}[Bounded Jacobian and lag budget]
\label{ass:app_bounded}
There exist finite constants $L,E$ such that for every round $t$,
\begin{equation}
    \|G_t\|_1
    \coloneqq
    \max_{j\in[K]}\sum_{k=1}^K |[G_t]_{kj}|
    \le L,
    \qquad
    0\le \epsilon_{t,k}\le E.
    \label{eq:app_boundedness}
\end{equation}
\end{assumption}

For a horizon $T$, define the realized dual magnitude
\begin{equation}
    \Lambda_T
    \coloneqq
    \max_{1\le t\le T+1}\|\lambda_t\|_\infty,
    \qquad
    H\coloneqq 2L+E.
    \label{eq:app_constants}
\end{equation}
Then Assumption~\ref{ass:app_bounded} implies
\begin{equation}
    \|s_t\|_\infty\le L(1+\Lambda_T),
    \qquad
    \|h_t(p_t)\|_\infty\le H.
    \label{eq:app_basic_bounds}
\end{equation}
The second inequality follows from $\|u-p_t\|_1\le 2$ and $|[G_t]_{kj}|\le L$.

\subsection{Constrained no-regret guarantee}
\label{app:theory_regret}

For a fixed comparator $p\in\Delta_K$, define cumulative utility regret
\begin{equation}
    R_T(p)
    \coloneqq
    \sum_{t=1}^T\bigl(F_t(p)-F_t(p_t)\bigr).
    \label{eq:app_regret}
\end{equation}
The next theorem gives the finite-horizon bound from which the $O(\sqrt{T})$ statement in the main text follows.

\begin{theorem}[Constrained no-regret of full-Jacobian PDCM]
\label{thm:app_no_regret}
Suppose Assumptions~\ref{ass:app_concavity}--\ref{ass:app_bounded} hold and $(p_t,\lambda_t)$ is generated by Equations~\eqref{eq:app_primal_update}--\eqref{eq:app_dual_update}. Then for every horizon $T\ge 1$, every $p\in\mathcal{P}_T$, and every $\eta_p,\eta_\lambda>0$,
\begin{equation}
    R_T(p)
    \le
    \frac{\log K}{\eta_p}
    +
    \frac{\eta_p T}{2}L^2(1+\Lambda_T)^2
    +
    \frac{\eta_\lambda T}{2}KH^2.
    \label{eq:app_finite_regret_bound}
\end{equation}
Moreover, for every task $k$,
\begin{equation}
    \left[
        \sum_{t=1}^T c_{t,k}(p_t)
    \right]_+
    \le
    \left[
        \sum_{t=1}^T h_{t,k}(p_t)
    \right]_+
    \le
    \frac{\Lambda_T}{\eta_\lambda}.
    \label{eq:app_constraint_bound}
\end{equation}
If the dual sequence is stable, i.e., $\Lambda_T\le\Lambda<\infty$ for all $T$, and $\eta_p,\eta_\lambda=\Theta(T^{-1/2})$, then
\begin{equation}
    R_T(p)=O(\sqrt{T}),
    \qquad
    \left[
        \sum_{t=1}^T c_{t,k}(p_t)
    \right]_+
    =O(\sqrt{T}),
    \label{eq:app_sqrt_rates}
\end{equation}
so both average regret and average excess counterfactual lag vanish as $O(T^{-1/2})$.
\end{theorem}

\begin{proof}
We first prove a one-step entropic mirror-ascent inequality. By optimality of Equation~\eqref{eq:app_primal_update} and the three-point identity for the negative-entropy Bregman divergence, for every $p\in\Delta_K$,
\begin{equation}
    \eta_p\langle s_t,p-p_{t+1}\rangle
    \le
    D_{\mathrm{KL}}(p\|p_t)
    -D_{\mathrm{KL}}(p\|p_{t+1})
    -D_{\mathrm{KL}}(p_{t+1}\|p_t).
    \label{eq:app_three_point}
\end{equation}
Adding $\eta_p\langle s_t,p_{t+1}-p_t\rangle$ to both sides and using H\"older's inequality together with Pinsker's inequality,
\begin{equation}
    \langle s_t,p_{t+1}-p_t\rangle
    \le
    \|s_t\|_\infty\|p_{t+1}-p_t\|_1,
    \qquad
    D_{\mathrm{KL}}(p_{t+1}\|p_t)
    \ge
    \frac{1}{2}\|p_{t+1}-p_t\|_1^2,
\end{equation}
we obtain
\begin{equation}
    \langle s_t,p-p_t\rangle
    \le
    \frac{
        D_{\mathrm{KL}}(p\|p_t)
        -D_{\mathrm{KL}}(p\|p_{t+1})
    }{\eta_p}
    +
    \frac{\eta_p}{2}\|s_t\|_\infty^2.
    \label{eq:app_mirror_one_step}
\end{equation}
Summing over $t$, using $p_1=u$, $D_{\mathrm{KL}}(p\|u)\le\log K$, and Equation~\eqref{eq:app_basic_bounds}, gives
\begin{equation}
    \sum_{t=1}^T\langle s_t,p-p_t\rangle
    \le
    \frac{\log K}{\eta_p}
    +
    \frac{\eta_p T}{2}L^2(1+\Lambda_T)^2.
    \label{eq:app_mirror_sum}
\end{equation}

We now relate the score to utility regret and the safety residuals. From Equation~\eqref{eq:app_affine_residual},
\begin{equation}
    h_t(p_t)-h_t(p)=G_t(p-p_t).
    \label{eq:app_h_difference}
\end{equation}
Therefore, using Equation~\eqref{eq:app_score},
\begin{align}
    \langle s_t,p-p_t\rangle
    &=
    \left\langle G_t^\top\mathbf{1},p-p_t\right\rangle
    +
    \left\langle \lambda_t,G_t(p-p_t)\right\rangle\\
    &=
    \left\langle G_t^\top\mathbf{1},p-p_t\right\rangle
    +
    \left\langle \lambda_t,h_t(p_t)-h_t(p)\right\rangle.
    \label{eq:app_score_decomposition}
\end{align}
For $p\in\mathcal{P}_T$, we have $h_t(p)\le0$ coordinatewise and $\lambda_t\ge0$, so
\begin{equation}
    \left\langle G_t^\top\mathbf{1},p-p_t\right\rangle
    +
    \langle\lambda_t,h_t(p_t)\rangle
    \le
    \langle s_t,p-p_t\rangle.
    \label{eq:app_safe_drop}
\end{equation}
Since each $U_{t,k}$ is concave, $F_t$ is concave, and Lemma~\ref{lem:app_tangent} implies
\begin{equation}
    F_t(p)-F_t(p_t)
    \le
    \langle\nabla F_t(p_t),p-p_t\rangle
    =
    \left\langle G_t^\top\mathbf{1},p-p_t\right\rangle.
    \label{eq:app_F_concavity}
\end{equation}
Combining Equations~\eqref{eq:app_mirror_sum}, \eqref{eq:app_safe_drop}, and \eqref{eq:app_F_concavity} yields
\begin{equation}
    R_T(p)
    +
    \sum_{t=1}^T\langle\lambda_t,h_t(p_t)\rangle
    \le
    \frac{\log K}{\eta_p}
    +
    \frac{\eta_p T}{2}L^2(1+\Lambda_T)^2.
    \label{eq:app_predual}
\end{equation}

It remains to control the dual term. Projection onto the nonnegative orthant is nonexpansive and fixes the origin, hence Equation~\eqref{eq:app_dual_update} gives
\begin{align}
    \|\lambda_{t+1}\|_2^2
    &\le
    \|\lambda_t+\eta_\lambda h_t(p_t)\|_2^2\\
    &=
    \|\lambda_t\|_2^2
    +2\eta_\lambda\langle\lambda_t,h_t(p_t)\rangle
    +\eta_\lambda^2\|h_t(p_t)\|_2^2.
    \label{eq:app_dual_square}
\end{align}
Rearranging and summing from $t=1$ to $T$, using $\lambda_1=0$ and $\|h_t(p_t)\|_2^2\le KH^2$, gives
\begin{equation}
    -\sum_{t=1}^T\langle\lambda_t,h_t(p_t)\rangle
    \le
    \frac{\eta_\lambda T}{2}KH^2.
    \label{eq:app_dual_inner_bound}
\end{equation}
Substituting Equation~\eqref{eq:app_dual_inner_bound} into Equation~\eqref{eq:app_predual} proves Equation~\eqref{eq:app_finite_regret_bound}.

For the constraint bound, the coordinatewise projection satisfies
\begin{equation}
    \lambda_{t+1,k}
    =
    \max\{0,\lambda_{t,k}+\eta_\lambda h_{t,k}(p_t)\}
    \ge
    \lambda_{t,k}+\eta_\lambda h_{t,k}(p_t).
\end{equation}
Telescoping and using $\lambda_{1,k}=0$ gives
\begin{equation}
    \eta_\lambda\sum_{t=1}^T h_{t,k}(p_t)
    \le
    \lambda_{T+1,k}
    \le
    \Lambda_T.
\end{equation}
Taking positive parts proves the second inequality in Equation~\eqref{eq:app_constraint_bound}. The first follows from Equation~\eqref{eq:app_true_vs_linearized} after summing over $t$. Finally, if $\Lambda_T\le\Lambda$ and both step sizes are of order $T^{-1/2}$, every term in Equations~\eqref{eq:app_finite_regret_bound} and \eqref{eq:app_constraint_bound} is $O(\sqrt{T})$.
\end{proof}

The preceding theorem immediately yields the comparator used in the main text.

\begin{corollary}[Best fixed feasible allocation]
\label{cor:app_best_fixed}
Let
\begin{equation}
    p_T^\star
    \in
    \arg\max_{p\in\mathcal{P}_T}
    \sum_{t=1}^T F_t(p).
    \label{eq:app_best_fixed}
\end{equation}
Under the dual-stability condition $\Lambda_T\le\Lambda$ and step sizes $\eta_p,\eta_\lambda=\Theta(T^{-1/2})$,
\begin{equation}
    \sum_{t=1}^T
    \left[F_t(p_T^\star)-F_t(p_t)\right]
    =O(\sqrt{T}),
    \label{eq:app_best_fixed_regret}
\end{equation}
and for every task $k$,
\begin{equation}
    \left[
        \sum_{t=1}^T
        \left(
            U_{t,k}(u)-U_{t,k}(p_t)-\epsilon_{t,k}
        \right)
    \right]_+
    =O(\sqrt{T}).
    \label{eq:app_best_fixed_constraint}
\end{equation}
\end{corollary}

\begin{corollary}[Uniform-mixture safeguard]
\label{cor:app_uniform}
Because $u\in\mathcal{P}_T$, Theorem~\ref{thm:app_no_regret} applies in particular to the uniform mixture. Under dual stability and $\eta_p,\eta_\lambda=\Theta(T^{-1/2})$,
\begin{equation}
    \sum_{t=1}^T F_t(u)
    -
    \sum_{t=1}^T F_t(p_t)
    =O(\sqrt{T}).
    \label{eq:app_uniform_regret}
\end{equation}
Hence the average aggregate-utility gap between PDCM and the static uniform mixture vanishes as $O(T^{-1/2})$. In addition,
\begin{equation}
    \frac{1}{T}\sum_{t=1}^T
    \left(U_{t,k}(u)-U_{t,k}(p_t)\right)
    \le
    \frac{1}{T}\sum_{t=1}^T\epsilon_{t,k}
    +O(T^{-1/2})
    \label{eq:app_uniform_lag}
\end{equation}
for every task $k$.
\end{corollary}

\begin{remark}[Dual stability]
\label{rem:app_dual_stability}
The finite-horizon bounds in Equations~\eqref{eq:app_finite_regret_bound}--\eqref{eq:app_constraint_bound} hold for the realized quantity $\Lambda_T$ and do not assume that the dual sequence is bounded in advance. The $O(\sqrt{T})$ rates require $\Lambda_T=O(1)$. This is a stability condition on the coupled allocation and training dynamics; it may be checked empirically, imposed as an assumption, or enforced by an upper projection on the dual variables when such a projection is inactive relative to the desired optimum.
\end{remark}

\subsection{Dual-weighted marginal-utility matching}
\label{app:theory_matching}

The no-regret result describes finite-horizon behavior. A complementary stationary characterization explains the allocation condition targeted by the primal-dual coupling. Let $U_k:\Delta_K\to\mathbb{R}$ be differentiable and concave, define
\begin{equation}
    U(p)
    \coloneqq
    \begin{bmatrix}
        U_1(p)&\cdots&U_K(p)
    \end{bmatrix}^{\!\top},
    \qquad
    G(p)\coloneqq\nabla_p U(p),
    \qquad
    F(p)\coloneqq\mathbf{1}^\top U(p),
\end{equation}
and consider the static constrained allocation problem
\begin{equation}
\begin{aligned}
    \max_{p\in\Delta_K}\quad
    &F(p)\\
    \text{subject to}\quad
    &U_k(u)-U_k(p)-\epsilon_k\le0,
    \qquad k\in[K].
\end{aligned}
\label{eq:app_stationary_problem}
\end{equation}
Since $F$ is concave and each constraint function $U_k(u)-U_k(p)-\epsilon_k$ is convex, Equation~\eqref{eq:app_stationary_problem} is a convex optimization problem after negating the objective.

\begin{theorem}[Dual-weighted marginal-utility matching]
\label{thm:app_matching}
Assume Equation~\eqref{eq:app_stationary_problem} satisfies Slater's condition, and let $(p^\star,\lambda^\star)$ be a primal-dual optimum. Define
\begin{equation}
    s^\star
    \coloneqq
    G(p^\star)^\top(\mathbf{1}+\lambda^\star).
    \label{eq:app_stationary_score}
\end{equation}
Then there exists $\mu^\star\in\mathbb{R}$ such that
\begin{equation}
    s_j^\star=\mu^\star
    \quad\text{for every }j\text{ with }p_j^\star>0,
    \qquad
    s_j^\star\le\mu^\star
    \quad\text{for every }j\text{ with }p_j^\star=0.
    \label{eq:app_matching_condition}
\end{equation}
The multipliers additionally satisfy
\begin{equation}
    \lambda_k^\star\ge0,
    \qquad
    U_k(u)-U_k(p^\star)-\epsilon_k\le0,
    \qquad
    \lambda_k^\star\bigl(U_k(u)-U_k(p^\star)-\epsilon_k\bigr)=0.
    \label{eq:app_complementarity}
\end{equation}
If $p_j^\star>0$ for every $j$, then $p^\star$ is a fixed point of the exponentiated-gradient primal map with $\lambda^\star$ if and only if $s^\star=\mu^\star\mathbf{1}$.
\end{theorem}

\begin{proof}
Rewrite Equation~\eqref{eq:app_stationary_problem} as the convex minimization problem
\begin{equation}
\begin{aligned}
    \min_{p\in\Delta_K}\quad
    &-F(p)\\
    \text{subject to}\quad
    &c_k(p)\coloneqq U_k(u)-U_k(p)-\epsilon_k\le0,
    \qquad k\in[K].
\end{aligned}
\end{equation}
Under Slater's condition, the Karush--Kuhn--Tucker conditions are necessary and sufficient. Introduce multipliers $\lambda_k\ge0$ for $c_k(p)\le0$, $\mu\in\mathbb{R}$ for $\mathbf{1}^\top p=1$, and $\alpha_j\ge0$ for $-p_j\le0$. The KKT Lagrangian is
\begin{equation}
    \mathcal{L}_{\mathrm{KKT}}(p,\lambda,\mu,\alpha)
    =
    -F(p)
    +\sum_{k=1}^K\lambda_k\bigl(U_k(u)-U_k(p)-\epsilon_k\bigr)
    +\mu(\mathbf{1}^\top p-1)
    -\alpha^\top p.
\end{equation}
Differentiating with respect to $p$ gives the stationarity condition
\begin{equation}
    -G(p^\star)^\top(\mathbf{1}+\lambda^\star)
    +\mu^\star\mathbf{1}
    -\alpha^\star
    =0.
    \label{eq:app_KKT_stationarity}
\end{equation}
Hence $s_j^\star=\mu^\star-\alpha_j^\star$. If $p_j^\star>0$, complementary slackness for $-p_j\le0$ gives $\alpha_j^\star=0$, so $s_j^\star=\mu^\star$. If $p_j^\star=0$, then $\alpha_j^\star\ge0$, so $s_j^\star\le\mu^\star$. The remaining KKT conditions give Equation~\eqref{eq:app_complementarity}.

For the fixed-point statement, suppose first that $p_j^\star>0$ for every $j$ and $s^\star=\mu^\star\mathbf{1}$. Then the exponentiated-gradient map gives
\begin{equation}
    \frac{
        p_j^\star\exp(\eta_p s_j^\star)
    }{
        \sum_{\ell=1}^K p_\ell^\star\exp(\eta_p s_\ell^\star)
    }
    =
    \frac{
        p_j^\star\exp(\eta_p\mu^\star)
    }{
        \exp(\eta_p\mu^\star)\sum_{\ell=1}^K p_\ell^\star
    }
    =p_j^\star.
\end{equation}
Conversely, if the map fixes an interior $p^\star$, division by $p_j^\star>0$ shows that $\exp(\eta_p s_j^\star)$ equals the same normalization constant for every $j$. Taking logarithms gives $s_j^\star=s_\ell^\star$ for all $j,\ell$, which is exactly $s^\star=\mu^\star\mathbf{1}$ for some scalar $\mu^\star$.
\end{proof}

Theorem~\ref{thm:app_matching} exposes the role of the full Jacobian. Expanding the stationary score coordinatewise,
\begin{equation}
    s_j^\star
    =
    \sum_{k=1}^K
    (1+\lambda_k^\star)
    \frac{\partial U_k(p^\star)}{\partial p_j}.
    \label{eq:app_cross_task_matching}
\end{equation}
Thus the dual variable of task $k$ amplifies the entire $k$th row of the allocation-response Jacobian, not only its diagonal entry. A training coordinate $j$ is valuable when increasing its allocation improves any evaluation task with high marginal utility, and its score receives additional weight when those tasks are under active lag pressure. At an interior stationary optimum, these dual-weighted cross-task marginal effects are equalized across all training coordinates.
